%% Manuscript for Journal of Power Sources
%% Elsevier elsarticle template
\documentclass[5p,times,number]{elsarticle}
\usepackage{amsmath,amssymb}
\usepackage{amsthm}
\usepackage[utf8]{inputenc}
\usepackage{booktabs}
\usepackage{multirow}
\usepackage{xcolor}
\usepackage{siunitx}
\usepackage{algorithm}
\usepackage{algorithmic}
\usepackage{graphicx}
\usepackage{stfloats}
\usepackage{enumitem}
\raggedbottom          
\renewcommand{\topfraction}{0.95}
\renewcommand{\bottomfraction}{0.95}
\renewcommand{\textfraction}{0.02}
\renewcommand{\floatpagefraction}{0.85}
\renewcommand{\dbltopfraction}{0.95}
\renewcommand{\dblfloatpagefraction}{0.85}
\setcounter{topnumber}{4}
\setcounter{bottomnumber}{4}
\setcounter{totalnumber}{8}
\setcounter{dbltopnumber}{4}
% Reduce spacing around floats
\setlength{\dbltextfloatsep}{6pt plus 2pt minus 2pt}
\setlength{\dblfloatsep}{6pt plus 2pt minus 2pt}
\setlength{\textfloatsep}{6pt plus 2pt minus 2pt}
\setlength{\intextsep}{6pt plus 2pt minus 2pt}
\setlength{\floatsep}{6pt plus 2pt minus 2pt}
\setlength{\abovecaptionskip}{4pt}
\setlength{\belowcaptionskip}{2pt}
\definecolor{supercolor}{RGB}{96,160,216}
\usepackage[colorlinks=true,linkcolor=supercolor,citecolor=supercolor,urlcolor=supercolor,filecolor=supercolor]{hyperref}


\makeatletter
\AtBeginDocument{%
  \hypersetup{linkcolor=supercolor,citecolor=supercolor,urlcolor=supercolor,filecolor=supercolor}%
  \def\NAT@open{\textcolor{supercolor}{[}}%
  \def\NAT@close{\textcolor{supercolor}{]}}%
  \def\NAT@sep{\textcolor{supercolor}{,}}%
  \def\NAT@cmt{\textcolor{supercolor}{,\space}}%
}
\makeatother


\providecommand{\doi}[1]{}

\usepackage{etoolbox}


\newcommand{\figref}[1]{\hyperref[#1]{\textup{Figure~\ref*{#1}}}}
\newcommand{\tabref}[1]{\hyperref[#1]{\textup{Table~\ref*{#1}}}}
\newcommand{\secref}[1]{\hyperref[#1]{\textup{Section~\ref*{#1}}}}
\newcommand{\ssref}[1]{\hyperref[#1]{\textup{\S\ref*{#1}}}}
\newcommand{\appref}[1]{\hyperref[#1]{\textup{Appendix~\ref*{#1}}}}
\newcommand{\propref}[1]{\hyperref[#1]{\textup{Proposition~\ref*{#1}}}}
\newcommand{\eqnref}[1]{\hyperref[#1]{\textup{Eq.~(\ref*{#1})}}}
\newcommand{\eqnrefp}[1]{\hyperref[#1]{\textup{Eq.~\ref*{#1}}}}
\newcommand{\eqnsref}[2]{\hyperref[#1]{\textup{Eqs.~(\ref*{#1})}}--\hyperref[#2]{\textup{(\ref*{#2})}}}
\newcommand{\refsref}[2]{\hyperref[#1]{\textup{Eqs.~\ref*{#1}}}--\hyperref[#2]{\textup{\ref*{#2}}}}

\graphicspath{{./}{figures/}}
\sisetup{
    separate-uncertainty = true,
    multi-part-units = single
}

\newtheorem{proposition}{Proposition}

\setlength{\emergencystretch}{2em}


\pdfstringdefDisableCommands{
    \let\fnmark\relax
    \let\cormark\relax
    \let\fnref\relax
    \let\corref\relax
    \def\cnotenum{}
    \let\@corref\relax
    \let\ead\relax
    \let\address\relax
    \def\figref#1{Figure~\ref{#1}}
    \def\tabref#1{Table~\ref{#1}}
    \def\secref#1{Section~\ref{#1}}
    \def\ssref#1{S\ref{#1}}
    \def\appref#1{Appendix~\ref{#1}}
    \def\propref#1{Proposition~\ref{#1}}
    \def\eqnref#1{Eq.~(\ref{#1})}
    \def\eqnrefp#1{Eq.~\ref{#1}}
    \def\eqnsref#1#2{Eqs.~(\ref{#1})--(\ref{#2})}
    \def\refsref#1#2{Eqs.~\ref{#1}--\ref{#2}}
}

\journal{Journal of Power Sources}

\begin{document}

\begin{frontmatter}

\title{A Modified Shepherd Electrochemical Model with Coupled Temperature--Aging Dynamics for Lithium-Ion Cells under Constant-Current and Random-Walk Discharge}

\author[scu-cs]{Yinuo Wang\fnref{fn1}}

\author[scu-phy]{Yike Yang\fnref{fn1}}

\author[scu-math]{Lizhi Wei\fnref{fn1}}

\author[scu-phy]{Zhiqiang Li\corref{cor1}}
\ead{zhiqiangli@scu.edu.cn}

\fntext[fn1]{These authors contributed equally to this work.}
\cortext[cor1]{Corresponding author.}
\address[scu-cs]{School of Computer Science, Sichuan University, Chengdu 610064, China}
\address[scu-phy]{College of Physics, Sichuan University, Chengdu 610064, China}
\address[scu-math]{College of Mathematics, Sichuan University, Chengdu 610064, China}

\begin{abstract}
This paper presents a modified Shepherd-based electrochemical modelling framework for lithium-ion cells operating under the heterogeneous discharge profiles characteristic of portable electronic systems.
A five-parameter Shepherd voltage formulation---retaining one-to-one physical correspondence between coefficients and electrochemical phenomena (open-circuit potential, mass-transport polarisation, ohmic loss, exponential transient zone)---is coupled in a unified continuous-time differential-algebraic system with a five-component load decomposition, a logistic capacity--temperature and Arrhenius resistance--temperature submodel, and a power-law cycle/$\sqrt{t}$ calendar aging submodel for time-to-empty (TTE) prediction.
Parameters identified from the XJTU lithium-ion battery dataset yield a voltage-tracking RMSE of \SI{17.85}{\milli\volt} under constant-current discharge, reducing error by over 55\% relative to the Th\'{e}venin-1RC model and by 75\% relative to a Nernst baseline.
Cross-batch validation transfers parameters from Batch-1 (2\,C) to Batch-2 (3\,C) with no re-fitting and a statistically non-significant RMSE gap, confirming C-rate generalisability.
Dynamic-load validation on the NASA Randomized Battery Usage Dataset---random-walk currents spanning 0.5--\SI{4}{\ampere} over 200 cycles---yields \SI{63.5}{\milli\volt} RMSE using the same constant-current parameters, approximately $3.2\times$ smaller than a Th\'{e}venin-1RC ECM identified on the same reference data, demonstrating that the Shepherd formulation transfers to time-varying loads without re-identification.
Ablation across a $3{\times}3$ temperature--aging grid quantifies the non-additive coupling between submodels: temperature coupling contributes up to 70\% TTE error and aging up to 55\%, with synergistic interaction terms not predictable from either submodel in isolation.
A receding-horizon model-predictive control strategy built on the Shepherd prediction core extends simulated runtime by up to 58\% with a computational cost of \SI{0.009}{\milli\second} per simulated second, enabling embedded deployment.
\end{abstract}

\begin{highlights}
\item Modified five-parameter Shepherd model achieves 17.85\,mV voltage RMSE on XJTU cells.
\item 55--75\% RMSE reduction over Th\'{e}venin-1RC, Th\'{e}venin-2RC, R$_{\mathrm{int}}$ and Nernst baselines.
\item Cross-batch transfer 2\,C$\to$3\,C with no re-fitting; bootstrap RMSE gap CI spans zero.
\item NASA random-walk validation: 63.5\,mV RMSE, $\sim 3.2\times$ better than CC-fitted Th\'{e}venin-1RC.
\item Hessian condition number $\kappa(H)\!\approx\!2.4\!\times\!10^{10}$ quantifies practical identifiability of the model.
\end{highlights}

\begin{keyword}
lithium-ion battery \sep Shepherd model \sep equivalent-circuit comparison \sep dynamic discharge validation \sep cross-batch transferability \sep temperature--aging coupling
\end{keyword}

\end{frontmatter}

% ============================================================
\section{Introduction}\label{sec:intro}
Lithium-ion cells in portable electronic systems operate under markedly different conditions from those of stationary storage or electric-vehicle batteries: discharge currents are highly dynamic and heterogeneous---screen brightness, system-on-chip (SoC)\footnote{Throughout this paper, ``SOC'' (all capitals) denotes state of charge, while ``SoC'' (in mixed case) refers to the system-on-chip hardware component.} utilisation, radio-frequency communication, GPS positioning, and background services jointly determine the instantaneous power draw, with the global portable-device fleet now exceeding six billion units~\cite{Benkhelifa2016,Carroll2010,Chen2014}.
Accurate prediction of the cell-level \emph{time-to-empty} (TTE)---the interval from the current state of charge (SOC) to a cutoff threshold---under such time-varying loads is a long-standing problem at the interface of electrochemical modelling and battery management~\cite{Benkhelifa2016,Liu2019ECM}.
For portable-cell applications, a useful battery model must simultaneously (i)~capture the terminal-voltage response with quantitatively interpretable electrochemical parameters, (ii)~remain valid across C-rates and across dynamic random-walk loads with the same parameter set, (iii)~track the temperature-dependent capacity and resistance variations that dominate cold-environment performance, and (iv)~account for cycle and calendar aging that progressively degrade both voltage and capacity over the cell's service life~\cite{Liu2019ECM,Liu2018ECM,Vetter2005}.

Existing battery modelling approaches span a wide spectrum.
At one extreme, physics-based electrochemical models (e.g., pseudo-two-dimensional models) offer mechanistic fidelity but demand extensive parameterisation and prohibitive computation for real-time deployment~\cite{Doyle1993,Newman2004,Mama2025}.
At the other, data-driven methods (neural networks, Gaussian processes) bypass physical insight and require large labelled datasets that are rarely available under diverse operating conditions~\cite{Chemali2018,Richardson2017,Xu2012ECM}.
Equivalent-circuit models (ECMs)---including Rint, first-order Th\'{e}venin, and second-order Th\'{e}venin networks---occupy a middle ground with modest computational cost, yet they parameterise voltage dynamics via abstract RC elements that lack direct electrochemical interpretability~\cite{Hu2012,You2018,He2012ECM}.
The classical Shepherd formulation~\cite{Shepherd1965} offers explicit terms for polarisation drop, ohmic loss, and initial transient behaviour with clear physical correspondence, but has been underexplored for portable-device TTE prediction that must simultaneously account for component-level power disaggregation, temperature coupling, and long-term aging.
Fuller~\cite{Fuller2014} demonstrated a Shepherd-type model for constant-power discharge with rate effects, yet without multi-component load coupling, temperature--aging dynamics, or energy management integration that modern portable systems demand.
Specifically, Fuller's model (i)~targets primary alkaline cells (non-rechargeable), whereas portable devices universally employ rechargeable lithium-ion chemistries with distinct voltage profiles and aging mechanisms; (ii)~incorporates Peukert-type rate derating for variable-current discharge but assumes isothermal conditions, omitting the Arrhenius-governed resistance increase and logistic capacity reduction that dominate lithium-ion performance variation across the $0$--$45\,^{\circ}$C operating range; (iii)~treats the load as a scalar constant power, without component-level disaggregation that would enable targeted subsystem throttling; and (iv)~does not model cycle or calendar aging, making it unsuitable for lifetime-aware energy management.
The present work retains Fuller's core advantage---physical transparency of the Shepherd voltage equation---while addressing all four limitations through the coupled temperature, aging, and component-level power submodels developed in \secref{sec:model}.

Several studies have characterised smartphone power consumption at the component level~\cite{Carroll2010,Chen2014,Murmuria2012}, revealing that screen, SoC, and radio subsystems dominate energy usage.
Recent work has extended these analyses to 5G and WiFi power profiling~\cite{Xu2020,Lauridsen2014,Sun2018,Negri2006}, but the coupling between load decomposition and electrochemical state evolution remains largely ad hoc.
On the aging front, capacity fade and resistance growth in lithium-ion cells have been modelled via power-law cycling degradation and Arrhenius-accelerated calendar aging~\cite{Vetter2005,Wang2011JPS,Schmalstieg2014JPS,Whittingham2012,Li2019,Kou2017}, while temperature effects on capacity and internal resistance are well characterised~\cite{Ali2020,Chen2025,He2017,Tang2019ECM}.
Liu et al.~\cite{Liu2019ECM} addressed battery charging management considering degradation costs, and Liu et al.~\cite{Liu2018ECM} optimised multi-objective charging patterns for lithium-ion cells, yet neither integrates real-time discharge-side TTE prediction under heterogeneous loads.

Despite the maturity of individual submodels, their joint integration into a unified continuous-time framework with experimentally validated cross-batch and dynamic-load transferability remains largely unexplored~\cite{Tian2019,Xiong2018,Mama2025,Wang2020RSER}.
Critically, this gap is not merely one of implementation convenience: the coupling itself produces \emph{non-additive} effects---for example, the interaction between temperature-reduced capacity and aging-induced resistance growth yields TTE deviations that cannot be predicted by either submodel in isolation (see the ablation study in \secref{sec:ablation}).
The main contributions of this paper are as follows.
\begin{enumerate}[leftmargin=*, label=\arabic*.]
    \item \textit{A modified Shepherd voltage formulation for lithium-ion cells under heterogeneous loads.}
    Building on the classical Shepherd equation~\cite{Shepherd1965} and the constant-power lithium-ion adaptation by Fuller~\cite{Fuller2014}, we retain the five-parameter electrochemically interpretable form---open-circuit potential, mass-transport polarisation ($K/\mathrm{SOC}$), ohmic loss, and exponential transient zone---and demonstrate that, after differential-evolution identification on the XJTU dataset~\cite{XJTU2020}, it achieves a voltage-tracking RMSE of \SI{17.85}{\milli\volt} under constant-current discharge, reducing error by 55--75\% relative to four standard equivalent-circuit baselines (Nernst-based, Rint, Th\'{e}venin-1RC, Th\'{e}venin-2RC) widely used in the lithium-ion battery management literature~\cite{Hu2012,Plett2004,You2018}.
    The 5-parameter count is competitive with Rint (9), 1RC (10), and 2RC (12) implementations.
    \item \textit{Cross-batch generalisation across C-rates without re-fitting.}
    Parameters identified on XJTU Batch-1 (2\,C, 8 cells) are applied directly to Batch-2 (3\,C, 15 cells); the bootstrap-estimated cross-batch RMSE gap is statistically non-significant (95\% CI: $[-3.88, +0.07]$\,mV, spanning zero).
    To our knowledge this transferability test---explicitly recommended in recent battery-modelling reviews~\cite{Xiong2018,Wang2020RSER}---has not previously been reported for Shepherd-type formulations on lithium-ion chemistries (\secref{sec:crossbatch}).
    \item \textit{Dynamic-load validation on the NASA Randomized Battery Usage Dataset, with head-to-head ECM benchmark.}
    Using the same constant-current parameters with no re-identification, the model is evaluated against random-walk discharge currents drawn uniformly from $[0.5, 4]$\,A and held in \SI{5}{\minute} segments, on four 18650 cells (RW3--RW6) over 200 cycles~\cite{Bole2014}, achieving an aggregate voltage RMSE of \SI{63.5}{\milli\volt}---approximately $3.2\times$ smaller than a Th\'{e}venin-1RC ECM identified under the same CC-fit$\to$RW-test protocol (\SI{200.0}{\milli\volt}). This provides the first quantitative head-to-head comparison of a Shepherd-type voltage core against a representative dynamic-capable ECM under truly time-varying loads, addressing a robustness concern that has limited the deployment of empirical voltage models in dynamic battery management contexts (\secref{sec:dynamic_val}).
    \item \textit{Coupled temperature--aging submodels with quantified non-additive interaction.}
    The Shepherd voltage core is embedded in a unified differential-algebraic system with a logistic capacity--temperature submodel, an Arrhenius resistance--temperature submodel, and power-law cycle / $\sqrt{t}$ calendar aging submodels~\cite{Vetter2005,Wang2011JPS,Schmalstieg2014JPS}.
    A systematic ablation across a $3\times3$ temperature--aging grid quantifies that the combined error of the coupled framework is not additive in the individual ablation errors, demonstrating interaction effects that justify the unified formulation (\secref{sec:ablation}).
    \item \textit{Engineering application: real-time TTE prediction and MPC energy management.}
    As a downstream application, we expose the validated Shepherd model as the prediction engine of a receding-horizon model-predictive controller for discharge-side power scheduling.
    With a computational cost of \SI{0.009}{\milli\second} per simulated second the controller is suitable for embedded BMS deployment and, in simulation, extends runtime by up to \SI{58}{\percent} relative to an unmanaged baseline (\secref{sec:energy_mgmt}).
    The modelling and validation contributions (1)--(4) stand independently of this controller study.
\end{enumerate}

The remainder of this paper is organised as follows.
\secref{sec:related} reviews related work.
\secref{sec:model} presents the unified model formulation.
\secref{sec:experiments} describes datasets, baselines, and evaluation protocols.
\secref{sec:results} reports results, discussion, and energy management implications.
\secref{sec:conclusion} concludes with limitations and future directions.

% ============================================================
\section{Related Work}\label{sec:related}
Battery modelling approaches span a wide spectrum of complexity and interpretability.
Physics-based electrochemical models (e.g., the pseudo-two-dimensional Doyle--Fuller--Newman formulation~\cite{Doyle1993,Newman2004}) offer mechanistic fidelity but require extensive parameterisation and prohibitive computation for real-time deployment~\cite{Mama2025}.
Data-driven methods (neural networks, Gaussian processes) bypass physical insight and demand large labelled datasets rarely available under diverse operating conditions~\cite{Chemali2018,Richardson2017,Xu2012ECM,Lipu2021}.
Equivalent-circuit models (ECMs)---including Rint, Th\'{e}venin-1RC, and 2RC networks---enjoy widespread use in battery management systems owing to their low computational cost; their parameters, however, are abstract circuit elements rather than electrochemical quantities, and online identification typically relies on Kalman-filter or recursive least-squares schemes that depend on the assumed model structure~\cite{Hu2012,Plett2004,You2018,He2012ECM,Birkl2015}.
The classical Shepherd formulation~\cite{Shepherd1965} decomposes terminal voltage into physically transparent terms (OCV, polarisation, ohmic drop, exponential transient), yet has been underexplored for lithium-ion cells operating under heterogeneous, time-varying discharge profiles.
Fuller~\cite{Fuller2014} demonstrated a Shepherd-type model for constant-power discharge of primary alkaline cells, incorporating Peukert-type rate effects to account for reduced effective capacity at high discharge currents.
While Fuller's model provides an instructive precedent for analytically tractable Shepherd-based runtime prediction, it assumes isothermal operation, treats the load as a scalar constant power without component decomposition, and does not model aging---assumptions that must be relaxed for rechargeable lithium-ion cells in portable devices.

On the load side, Carroll and Heiser~\cite{Carroll2010} and Chen et al.~\cite{Chen2014} provided comprehensive smartphone power breakdowns, identifying screen, SoC, and radio as dominant consumers.
Subsequent work characterised LTE~\cite{Lauridsen2014}, 5G~\cite{Xu2020}, WiFi~\cite{Sun2018}, and Bluetooth~\cite{Negri2006} profiles.
However, these studies report measurements in isolation, without coupling to electrochemical state evolution.
On the degradation front, temperature effects on capacity and resistance are well characterised via logistic and Arrhenius models~\cite{Ali2020,Chen2025,He2017,Whittingham2012,Tang2019ECM}, while cycle aging follows power-law relations~\cite{Vetter2005,Wang2011JPS,Kou2017,Li2019} and calendar aging exhibits $\sqrt{t}$ dependence on time and Arrhenius dependence on temperature~\cite{Vetter2005,Schmalstieg2014JPS}.
Liu et al.~\cite{Liu2019ECM} proposed a coupled electrothermal-aging model for economy-conscious battery charging management that jointly captures electrical, thermal, and degradation dynamics across different timescales using a second-order RC circuit, a lumped thermal model, and a power-law capacity-loss submodel.
Their multi-timescale coupling architecture---fast electrical dynamics within a single charging cycle, slow aging dynamics across cycles---is a direct methodological precursor to the present work.
However, Liu et al.'s framework is designed for the \emph{charging} side in \emph{electric-vehicle} applications, and it does not address (a)~discharge-side TTE prediction, (b)~component-level load decomposition for heterogeneous portable-device workloads, or (c)~the use of electrochemically interpretable voltage models (Shepherd) in place of abstract RC circuits.
Liu et al.~\cite{Liu2018ECM} further optimised multi-objective charging patterns for lithium-ion cells.
The present work extends this coupled-submodel philosophy from \emph{charging optimisation} to \emph{discharge-side TTE prediction}, from \emph{EV-scale} to \emph{portable-device-scale} loads, and from \emph{ECM-based} to \emph{Shepherd-based} voltage dynamics with explicit electrochemical interpretability.
Three specific technical advances beyond Liu et al.~\cite{Liu2019ECM} merit emphasis: (i)~the five-component power disaggregation enables \emph{targeted} load throttling (reducing a specific subsystem rather than blanket power reduction), (ii)~the Shepherd voltage core provides physically transparent parameter--phenomenon correspondence absent in RC-circuit representations, and (iii)~the cross-batch validation (2\,C$\to$3\,C, no re-fitting) directly tests and confirms the model's transferability across C-rates---a robustness test not reported in the prior work.
Despite the maturity of individual submodels, their joint integration into a portable-device TTE framework with experimentally validated cross-condition generalisation has not been demonstrated~\cite{Tian2019,Xiong2018,Mama2025,Yu2025}.

\tabref{tab:gap} summarises the capability coverage of representative prior works.
\begin{table}[!tb]
\caption{Capability comparison of representative battery modelling approaches.}\label{tab:gap}
\centering
\resizebox{\columnwidth}{!}{
\begin{tabular}{@{}lccccc@{}}
\toprule
Approach & Interp.\ V & Load & Temp. & Aging & Mgmt \\
\midrule
P2D~\cite{Doyle1993} & \checkmark & -- & \checkmark & -- & -- \\
Shepherd (Fuller)~\cite{Fuller2014} & \checkmark & -- & -- & -- & -- \\
ECMs~\cite{Hu2012,You2018} & -- & -- & partial & -- & -- \\
ECM + aging (charge)~\cite{Liu2019ECM} & -- & -- & \checkmark & \checkmark & charge \\
ECM + uncertainty~\cite{Tang2019ECM} & -- & -- & \checkmark & \checkmark & -- \\
Data-driven~\cite{Chemali2018,Lipu2021} & -- & -- & implicit & implicit & -- \\
Profiling~\cite{Carroll2010} & -- & \checkmark & -- & -- & -- \\
\textbf{This work} & \checkmark & \checkmark & \checkmark & \checkmark & \textbf{discharge} \\
\bottomrule
\end{tabular}
}
\end{table}

% ============================================================
\section{Proposed Unified Continuous-Time Model}\label{sec:model}
This section develops the proposed framework from first principles.
The model comprises four coupled submodels---terminal-voltage dynamics (\ssref{sec:shepherd}), component-level power decomposition (\ssref{sec:power}), temperature--parameter coupling (\ssref{sec:temperature}), and aging dynamics (\ssref{sec:aging})---which are assembled into a unified differential-algebraic system (\ssref{sec:ode}) and solved numerically (\ssref{sec:numerics}).
\figref{fig:framework} illustrates the information flow among submodels.
\begin{figure}[!tb]
\centering
\includegraphics[width=\columnwidth]{fig_framework.pdf}
\caption{Architecture of the proposed unified framework showing information flow among the four coupled submodels and the TTE computation module.}\label{fig:framework}
\end{figure}

\subsection{Terminal-voltage submodel}\label{sec:shepherd}
\subsubsection{Physical basis of the Shepherd formulation}
The terminal voltage of a lithium-ion cell during discharge is governed by three concurrent phenomena~\cite{Shepherd1965}:
\begin{enumerate}[leftmargin=*, label=\arabic*.]
    \item \emph{Thermodynamic equilibrium.} The open-circuit voltage~$E_0$ reflects the Gibbs free-energy difference between the electrode active materials at full charge.
    \item \emph{Mass-transport polarisation.} As discharge proceeds, lithium-ion depletion near the electrode surface creates a concentration gradient.  The resulting overpotential is proportional to the current~$I$ and inversely proportional to the remaining charge fraction~SOC, yielding a term~$K I/\mathrm{SOC}$.  Physically, $K$ ($\Omega$) encapsulates the effective diffusion resistance of the intercalation host.
    \item \emph{Ohmic loss and initial transient.}  Series resistance~$R$ produces an instantaneous $IR$ drop, while the exponential zone ($A$, $B$) captures the initial voltage relaxation associated with double-layer charging and surface-film effects during early discharge.
\end{enumerate}
Combining these contributions, the \emph{modified Shepherd equation} for terminal voltage reads
\begin{equation}\label{eq:shepherd}
V_{\mathrm{term}} = E_0 - \frac{K}{\mathrm{SOC}}\,I - R\,I + A\exp\!\bigl(-B Q_0(1-\mathrm{SOC})\bigr),
\end{equation}
where the four right-hand-side terms represent, from left to right, the open-circuit voltage, polarisation overpotential, ohmic drop, and exponential transient zone.
In the taxonomy of battery-scale electrical models surveyed by Mama et al.~\cite{Mama2025}, the Shepherd formulation occupies a distinct niche between purely empirical models and full equivalent-circuit networks: its coefficients are empirically determined---readily obtainable from standard discharge curves~\cite{Fuller2014}---yet each term retains a one-to-one mapping to an identifiable electrochemical phenomenon.
This combination of low parameterisation cost and physical transparency motivates its adoption as the voltage core of the present framework.

\subsubsection{SOC dynamics via Coulomb counting}
Conservation of charge relates the SOC evolution to the discharge current:
\begin{equation}\label{eq:coulomb}
\frac{\mathrm{d}\,\mathrm{SOC}}{\mathrm{d}t} = -\frac{I(t)}{3600\,Q_0},
\end{equation}
where $Q_0$ (Ah) is the effective cell capacity and the factor 3600 converts Ah to A$\cdot$s.  The sign convention is $I>0$ during discharge.

\subsubsection{Power-balance closure}
For a smartphone whose instantaneous total power consumption is $P_{\mathrm{total}}(t)$, the algebraic constraint
\begin{equation}\label{eq:power}
V_{\mathrm{term}}(\mathrm{SOC}, I) \cdot I = P_{\mathrm{total}}(t)
\end{equation}
closes the system by implicitly determining $I$ at each instant.

\subsubsection{Time-to-empty definition}
Given an initial condition $\mathrm{SOC}(0) = \mathrm{SOC}_0$, the TTE is
\begin{equation}\label{eq:tte}
\mathrm{TTE} = \inf\bigl\{t > 0 \;\big|\; \mathrm{SOC}(t) \le \mathrm{SOC}_{\min} \;\text{or}\; V_{\mathrm{term}}(t) \le V_{\mathrm{cutoff}}\bigr\},
\end{equation}
where $\mathrm{SOC}_{\min} = 0.05$ and $V_{\mathrm{cutoff}} = \SI{2.5}{\volt}$ in this study.

\paragraph{Comparison with ECM representations}
In a Th\'{e}venin-$n$RC model, the $n$ RC-voltage state variables evolve via linear ODEs whose time constants ($\tau_k = R_k C_k$) lack one-to-one correspondence with electrochemical processes.
The Shepherd form, by contrast, replaces these abstract dynamics with a single, physically transparent $K/\mathrm{SOC}$ polarisation term and an exponential transient zone, reducing the parameter count from $2n+1$ (Th\'{e}venin) to~5 while yielding superior accuracy on constant-rate discharge profiles (see \secref{sec:voltage}).

\subsection{Component-level power decomposition}\label{sec:power}
A smartphone's instantaneous power consumption is disaggregated into five hardware subsystems~\cite{Carroll2010,Chen2014}:
\begin{equation}\label{eq:ptotal}
P_{\mathrm{total}} = P_{\mathrm{screen}} + P_{\mathrm{SoC}} + P_{\mathrm{radio}} + P_{\mathrm{GPS}} + P_{\mathrm{base}}.
\end{equation}
\paragraph{Screen power}
OLED display power depends on luminance~$L$ (cd/m$^2$), active area~$A_s$ (m$^2$), and a refresh-rate scaling factor~$\eta$:
\begin{equation}\label{eq:screen}
P_{\mathrm{screen}} = \eta\bigl(\gamma_s A_s L + \lambda_s A_s\bigr),
\end{equation}
where $\gamma_s$ (W\,cd$^{-1}$) and $\lambda_s$ (W/m$^{2}$) are OLED-specific power coefficients calibrated from~\cite{Lu2008}, and
\begin{equation}
\eta = \begin{cases} 0 & \text{screen off,} \\ 1.0 & 60\,\text{Hz refresh,} \\ 1.3 & 120\,\text{Hz refresh.} \end{cases}
\end{equation}
The factor $\eta = 1.3$ for 120\,Hz reflects the roughly 30\% increase in display driver power measured in recent high-refresh-rate panels.

\paragraph{Computational, communication, and baseline subsystems}
Each remaining component's power is parameterised by a dimensionless \emph{working intensity}~$\xi_i \in [0,1]$, representing the fraction of its maximum rated power:
\begin{equation}\label{eq:xi}
\begin{pmatrix} P_{\mathrm{SoC}} \\ P_{\mathrm{radio}} \\ P_{\mathrm{base}} \end{pmatrix}
= \operatorname{diag}(\xi_1,\,\xi_2,\,\xi_3)
\begin{pmatrix} P_{\mathrm{SoC}}^{\max} \\ P_{\mathrm{radio}}^{\max} \\ P_{\mathrm{base}}^{\max} \end{pmatrix}.
\end{equation}
GPS power is modelled as binary: $P_{\mathrm{GPS}} = P_{\mathrm{GPS}}^{\max}$ when satellite acquisition is active and zero otherwise, reflecting the all-or-nothing nature of positioning hardware.
Reference maximum powers adopted in this study are $P_{\mathrm{SoC}}^{\max} = \SI{2.04}{\watt}$, $P_{\mathrm{radio}}^{\max} = \SI{1.173}{\watt}$, $P_{\mathrm{GPS}}^{\max} = \SI{0.30}{\watt}$, and $P_{\mathrm{base}}^{\max} = \SI{0.12}{\watt}$, consistent with measurements reported in~\cite{Carroll2010,Chen2014}.
The SoC utilisation values in \tabref{tab:scenarios} include a 10\% baseline allocation for OS housekeeping processes.

\paragraph{Extensibility}
This additive decomposition is modular: adding a new subsystem (e.g., NFC, UWB) amounts to appending a $(\xi_j, P_j^{\max})$ pair.

\subsection{Temperature--parameter coupling}\label{sec:temperature}
Temperature influences cell behaviour through three mechanisms, each modelled with a distinct functional form.
\subsubsection{Capacity--temperature response}
The maximum number of lithium ions accommodated in the electrode lattice follows a Fermi--Dirac (logistic) occupancy distribution.
Denoting by $Q_0'$ the limiting capacity at infinite temperature, the effective capacity at temperature~$T$ (in Kelvin) is
\begin{equation}\label{eq:Q0T}
Q_0^{\mathrm{th}}(T) = \frac{Q_0'}{1 + \exp\!\Bigl(\alpha_\theta\bigl(\tfrac{1}{T} - \tfrac{1}{T_m}\bigr)\Bigr)},
\end{equation}
where $\alpha_\theta$ (K) controls the steepness and $T_m$ (K) is the midpoint temperature.
Fitted values (see \secref{sec:paramfit}) are $\alpha_\theta = \SI{3000.0}{\kelvin}$ and $T_m = \SI{298.51}{\kelvin}$.
\emph{Physical interpretation.}
At low temperatures ($T \ll T_m$), the exponential term dominates, sharply reducing $Q_0^{\mathrm{th}}$.
This reflects the thermodynamic freezing-out of available lattice sites, consistent with experimental observations of reduced deliverable capacity below \SI{0}{\degreeCelsius}~\cite{Ali2020}.

\subsubsection{Resistance--temperature response}
The temperature dependence of the total internal resistance is governed by the Arrhenius law for thermally activated charge-transfer and ion-transport processes:
\begin{equation}\label{eq:RT}
R^{\mathrm{th}}(T) = R_0\exp\!\biggl(\frac{E_r}{k_B}\Bigl(\frac{1}{T} - \frac{1}{T_0}\Bigr)\biggr),
\end{equation}
where $E_r = \SI{0.005}{\electronvolt}$ is the effective activation energy, $k_B$ is the Boltzmann constant, and $T_0 = \SI{298.15}{\kelvin}$ is the reference temperature~\cite{Whittingham2012,He2017,Jeon2011ECM}.
$R_0$ denotes the resistance at $T_0$.
The small value of $E_r$ indicates that temperature-induced resistance variation is modest ($<2\%$ over the operating range), with TTE sensitivity to temperature dominated instead by the logistic capacity response (\eqnrefp{eq:Q0T}).

\subsubsection{SoC leakage power}
At elevated temperatures, subthreshold leakage current in the SoC increases according to a MOSFET leakage model~\cite{Bhat2019,Bhat2021}:
\begin{equation}\label{eq:leak}
P_{\mathrm{leak}}(T) = V_{\mathrm{nom}}\,c_1\,T^2\exp(c_2 / T),
\end{equation}
with $V_{\mathrm{nom}} = \SI{3.7}{\volt}$ (the nominal cell voltage), $c_1 = \SI{0.493}{\ampere\per\kelvin\squared}$, and $c_2 = \SI{-3740}{\kelvin}$.
This term is additive to the SoC dynamic power in \eqnref{eq:ptotal}.

\emph{Scope.}
Because $P_{\mathrm{leak}}(T_{\mathrm{ref}}) \approx 0$ at the \SI{25}{\degreeCelsius} reference, leakage has negligible impact on the nominal TTE results in \secref{sec:tte}.
It is therefore included only in the temperature-sweep analysis (\secref{sec:temp_tte}) where the high-temperature regime is explicitly explored, and omitted from the baseline TTE and ablation simulations to avoid introducing a temperature-dependent power offset that would obscure the comparison of electrochemical submodel contributions.

\subsection{Aging dynamics}\label{sec:aging}
Repeated charge--discharge cycling degrades cell performance through SEI layer growth, lithium plating, and active-material loss~\cite{Vetter2005}.
We model both \emph{cycle aging} and \emph{calendar aging}.
\subsubsection{Cycle aging: capacity fade and resistance growth}
Denoting by $N$ the number of equivalent full cycles, cycle-induced degradation is captured by power-law relations fitted to XJTU data:
\paragraph{Resistance growth}
\begin{equation}\label{eq:Raging}
R^{\mathrm{cyc}}(N) = R_0^{\mathrm{nom}}\bigl(1 + \beta_R\, N^{\gamma_R}\bigr),
\end{equation}
with $\beta_R = 1.229 \times 10^{-3}$ and $\gamma_R = 1.085$.  This super-linear growth ($\gamma_R > 1$) reflects SEI thickening and progressive pore clogging.
\paragraph{Capacity fade}
\begin{equation}\label{eq:Qaging}
Q_0^{\mathrm{cyc}}(N) = Q_0^{\mathrm{nom}}\bigl(1 - \alpha_Q\, N^{\gamma_Q}\bigr),
\end{equation}
with $\alpha_Q = 2.56 \times 10^{-4}$ and $\gamma_Q = \gamma_R$.  Using a shared exponent $\gamma$ reflects the common mechanistic origin (SEI growth) underlying both degradation modes.

\subsubsection{Calendar aging}
During storage, diffusion-limited SEI growth produces a $\sqrt{t_s}$ capacity loss:
\begin{equation}\label{eq:calaging}
\Delta Q_0^{\mathrm{cal}}(t_s, T, \mathrm{SOC}) = Q_0^{\mathrm{nom}}\,\kappa\,\sqrt{t_s}\,\exp\!\Bigl(-\frac{E_a}{R_g T} + \lambda_a\,\mathrm{SOC}\Bigr),
\end{equation}
where $\kappa = 10^{-5}$ (s$^{-1/2}$) is the calendar aging coefficient, $E_a = \SI{50}{\kilo\joule\per\mole}$ is the activation energy for SEI growth, $R_g = \SI{8.314}{\joule\per\mole\per\kelvin}$ is the universal gas constant, and $\lambda_a = 1.5$ is the SOC-stress coefficient capturing the dependence on anode potential (higher SOC $\Rightarrow$ higher electrode potential $\Rightarrow$ faster parasitic reactions). The numerical values for $\kappa$, $E_a$, and $\lambda_a$ are adopted as representative orders of magnitude for NMC/graphite chemistries from the holistic NMC aging study of Schmalstieg et al.~\cite{Schmalstieg2014JPS} and the comprehensive degradation review of Vetter et al.~\cite{Vetter2005}; recalibration on cycle-resolved storage data is recommended for cells of different chemistry or format.

\subsubsection{Combined aging--temperature coupling}\label{sec:coupling}
The effective capacity and resistance under simultaneous temperature and aging effects are obtained by composing the submodel outputs:
\begin{align}
Q_0(T, N, t_s) &= Q_0^{\mathrm{th}}(T) \cdot \frac{Q_0^{\mathrm{cyc}}(N) - \Delta Q_0^{\mathrm{cal}}}{Q_0^{\mathrm{nom}}}, \label{eq:Q0coupled}\\[4pt]
R(T, N) &= R^{\mathrm{th}}(T) \cdot \frac{R^{\mathrm{cyc}}(N)}{R_0^{\mathrm{nom}}}. \label{eq:Rcoupled}
\end{align}
\emph{Rationale.}
The multiplicative coupling in \eqnsref{eq:Q0coupled}{eq:Rcoupled} assumes that temperature and aging act on \emph{independent} physical mechanisms---thermal reduction of available lattice sites (\eqnrefp{eq:Q0T}) vs.\ electrochemical loss of active material (\refsref{eq:Qaging}{eq:calaging})---so their effects compose as scaling factors on the nominal values.
This separability assumption is common in the aging literature~\cite{Vetter2005,Li2019,Tang2019ECM} and is supported by the observation that temperature-dependent capacity tests at different cycle counts yield approximately parallel logistic curves (see XJTU temperature-compensation data).

\subsection{Coupled ODE system and TTE computation}\label{sec:ode}
Assembling the submodels, the complete state-evolution system is:
\begin{equation}\label{eq:system}
\left\{\;
\begin{aligned}
& V_{\mathrm{term}} = E_0 - \frac{K}{\mathrm{SOC}}\,I - R(T,N)\,I + A\,\mathrm{e}^{-B\,Q_0(T,N)(1-\mathrm{SOC})}\\[6pt]
& \frac{\mathrm{d}\,\mathrm{SOC}}{\mathrm{d}t} = -\frac{I(t)}{3600\,Q_0(T,N)}\\[6pt]
& V_{\mathrm{term}}(I) \cdot I = P_{\mathrm{total}}(t)\\[4pt]
& \mathrm{SOC}(0) = \mathrm{SOC}_0
\end{aligned}
\right.
\end{equation}
This is a differential-algebraic equation (DAE) of index~1: the algebraic constraint (third line) can be solved for $I$ at each time step, after which the SOC ODE is integrated forward.
\paragraph{State vector}
The minimal state is $\mathbf{x}(t) = [\mathrm{SOC}(t)]$ (one-dimensional), since temperature~$T$ and cycle count~$N$ are treated as slowly varying external parameters (quasi-static assumption).
\paragraph{Termination}
Integration halts at $t^* = \mathrm{TTE}$ when $\mathrm{SOC}(t^*) \le 0.05$ or $V_{\mathrm{term}}(t^*) \le V_{\mathrm{cutoff}}$.

\subsection{Parameter identification}\label{sec:paramfit}
\subsubsection{Optimisation procedure}
The six Shepherd parameters $\boldsymbol{\theta} = (E_0, K, R_0, A, B, Q_0^{\mathrm{nom}})$ are identified by minimising the voltage RMSE over the first discharge cycle of each Batch-1 cell:
\begin{equation}\label{eq:obj}
\boldsymbol{\theta}^* = \arg\min_{\boldsymbol{\theta}} \sqrt{\frac{1}{M}\sum_{k=1}^{M}\bigl(V_{\mathrm{term}}^{\mathrm{pred}}(t_k;\boldsymbol{\theta}) - V_{\mathrm{meas}}(t_k)\bigr)^2},
\end{equation}
where $M$ is the number of time samples per discharge trajectory.
We use the \emph{differential evolution} (DE) global optimiser (SciPy, strategy ``best1bin'', population size 50, 1000 generations) with bound constraints reflecting physical plausibility (e.g., $E_0 \in [3.0, 4.5]$\,V, $R_0 \in [0.01, 0.2]$\,$\Omega$).

\subsubsection{Structural identifiability}
The Shepherd model has five independent algebraic terms, each involving a distinct functional dependence on $(I, \mathrm{SOC})$:
(i)~a constant ($E_0$),
(ii)~a term linear in $I$ ($R_0$),
(iii)~a term proportional to $I/\mathrm{SOC}$ ($K$),
(iv)~an exponential in $(1-\mathrm{SOC})$ ($A$, $B$).
Since a full discharge trajectory provides data spanning $\mathrm{SOC} \in [0.05, 1]$ with known $I(t)$, these four functional forms are linearly independent, ensuring \emph{structural identifiability} of all five parameters.
In practice, we verify \emph{practical identifiability} along two complementary axes.
First, 20 independent DE optimisation trials with different random seeds yield parameter estimates whose coefficient of variation is below 2\%, confirming that the global optimum is unique and robustly located.
Second, we quantify the local curvature of the cost surface at the optimum by computing the (scaled) Hessian $H = \nabla_{\boldsymbol{\theta}}^{2}\,J(\boldsymbol{\theta}^{*})$ via central finite differences (relative step $h_{\mathrm{rel}} = 2\times 10^{-3}$) on the pooled XJTU Batch-1 cost function (8 cells, $\sim$140k samples). The resulting eigenvalue spectrum, $\{1.0\!\times\!10^{-9},\, 2.9\!\times\!10^{-3},\, 1.6\!\times\!10^{-2},\, 8.5\!\times\!10^{-2},\, 24.9\}$, gives a condition number $\kappa(H) \approx 2.4 \times 10^{10}$. Inspection of the lowest-eigenvalue eigenvector reveals an almost-flat direction in the $(E_0, R_0)$ subspace---a structural co-linearity expected under constant-current loading, where the constant offset $E_0$ and the ohmic drop $R_0\,\bar{I}$ are partly interchangeable. The remaining four eigen-directions are well excited, and the four polarisation--exponential parameters $(K, A, B)$ together with the SOC-dependent coupling are identified with relative standard errors of $11\%$--$30\%$. The (E$_0$, R$_0$) ambiguity is resolved in practice by either (i) constraining $R_0$ from a separate pulse test or (ii) including the random-walk current segments of \secref{sec:dynamic_val} in the fit, both of which break the (E$_0, R_0$) co-linearity by exciting non-equilibrium ohmic transients.

\subsubsection{Aging-parameter estimation}
The aging coefficients $(\alpha_Q, \beta_R, \gamma)$ are fitted separately to the cycle-resolved capacity and resistance trajectories extracted from the XJTU dataset (Batch-1), using nonlinear least squares on \eqnsref{eq:Raging}{eq:Qaging}.

\subsection{Numerical solution}\label{sec:numerics}
\subsubsection{Algebraic current solver}
At each time step, the nonlinear algebraic equation
\begin{equation}\label{eq:rootfind}
f(I) \equiv V_{\mathrm{term}}(I, \mathrm{SOC}, T, N) \cdot I - P_{\mathrm{total}} = 0
\end{equation}
is solved for $I \ge 0$.
\begin{proposition}[Uniqueness of the physical root]\label{prop:mono}
Under the operating regime $E_0, K, R, A, B, Q_0 > 0$ and $\mathrm{SOC} \in (0, 1]$, let $F = E_0 + A\exp(-BQ_0(1-\mathrm{SOC}))$ and $\alpha = K/\mathrm{SOC} + R$.
The power function $g(I) = I\,V_{\mathrm{term}}(I) = FI - \alpha I^2$ is a concave quadratic with a unique maximum $g^* = F^2/(4\alpha)$.
For any $P_{\mathrm{total}} \in (0, g^*]$, \eqnref{eq:rootfind} has exactly one root on $(0, F/(2\alpha)]$, corresponding to the high-voltage (low-current) operating point.
\end{proposition}
The proof is given in \appref{app:proof}.
We implement a relaxed fixed-point iteration.
First, the initial guess is set as $I^{(0)} = P_{\mathrm{total}}/E_0$.
Second, the current is updated via $\tilde{I}^{(k+1)} = P_{\mathrm{total}} / V_{\mathrm{term}}(I^{(k)})$, followed by exponential-moving-average (EMA) smoothing $I^{(k+1)} = \rho\, I^{(k)} + (1-\rho)\,\tilde{I}^{(k+1)}$ with $\rho = 0.9$.
Finally, under Proposition~\ref{prop:mono}'s monotonicity, the EMA iteration contracts to the unique physical root within $\sim$10 steps for all operating points encountered in this study.
Because the Shepherd voltage is a simple quadratic in $I$ (Proposition~\ref{prop:mono}), this lightweight scheme achieves the same converged current as a Newton--Raphson solver while avoiding derivative evaluation, and is well suited to embedded deployment.

\subsubsection{Time integration}
SOC is advanced by explicit Euler:
\begin{equation}
\mathrm{SOC}_{n+1} = \mathrm{SOC}_n - \frac{I_n\,\Delta t}{3600\,Q_0},
\end{equation}
with a fixed step $\Delta t = \SI{1}{\second}$ chosen for reproducibility; the factor 3600 converts $Q_0$ from Ah to A$\cdot$s.
For stiff thermal--electrical couplings in future extensions, adaptive Runge--Kutta or implicit solvers (e.g., CVODE) are recommended.

\subsubsection{Computational cost}\label{sec:comp_cost}
Each time step requires one EMA fixed-point update plus one Euler step.
The proposed model simulates a full 6.4\,h discharge in \SI{0.20}{\second} (wall clock), corresponding to \SI{0.009}{\milli\second} per simulated second---approximately four orders of magnitude faster than finite-element P2D solvers~\cite{Doyle1993,Newman2004,Mama2025} and comparable to compiled ECM implementations, enabling real-time deployment on embedded hardware~\cite{Xiong2018,Wang2020RSER}.

% ============================================================
\section{Experimental Setup}\label{sec:experiments}
\subsection{Datasets}\label{sec:datasets}
\paragraph{XJTU Battery Dataset}
We use the publicly available XJTU lithium-ion battery dataset~\cite{XJTU2020}, comprising six batches of NMC/graphite 18650 cells.
Two batches are employed in this study.
Specifically, Batch-1 comprises 8 cells cycled at 2\,C constant-current discharge, serving as the \emph{fitting set} for parameter identification, while Batch-2 comprises 15 cells cycled at 3\,C constant-current discharge, used \emph{exclusively} for cross-batch generalisation testing.
Each cell record contains time-stamped voltage, current, temperature, and capacity measurements across hundreds of charge--discharge cycles.

\paragraph{Synthetic smartphone load profiles}
The XJTU dataset provides cell-level electrochemical data under controlled cycling conditions (constant current) on a battery test bench.
No real-time smartphone hardware telemetry (e.g., concurrent screen brightness, CPU frequency, radio state) is available.
Therefore, the component-level power profiles used in the TTE simulations (\tabref{tab:scenarios}) are \emph{synthetic}: they are constructed from published per-component power measurements on real smartphone hardware~\cite{Carroll2010,Chen2014,Lu2008,Murmuria2012} and applied as inputs to the electrochemically parameterised model.
This approach decouples voltage-model validation (performed on real cell data) from load-profile specification (constructed from literature), enabling systematic exploration of usage scenarios that would be infeasible with purely empirical data collection.
The implications and limitations of this separation are discussed in \secref{sec:conclusion}.

\paragraph{NASA Randomized Battery Usage Dataset}
To validate the Shepherd model under dynamic (non-constant-current) discharge conditions, we employ Dataset~\#11 from the NASA Prognostics Center of Excellence~\cite{Bole2014}.
This dataset comprises four 18650 lithium-ion cells (RW3--RW6, $\sim$\SI{2}{\ampere\hour}) subjected to two types of discharge:
(i)~\emph{reference CC discharges} at $\sim$\SI{1}{\ampere} to a \SI{3.2}{\volt} cutoff, recorded periodically as capacity benchmarks; and
(ii)~\emph{random-walk discharges}, where the current is drawn from a uniform distribution $I \sim \mathrm{U}(0.5, 4)$\,A and held constant for \SI{5}{\minute} segments until the cutoff is reached.
Each cell contains $\sim$44 reference discharges and $\sim$50 random-walk cycles.
This dataset is uniquely suited for testing the Shepherd model's transferability from CC-identified parameters to arbitrary current profiles, as the random-walk protocol spans a wide C-rate range ($0.25$--$2$C) within a single discharge.

\paragraph{Data preprocessing}
Discharge segments are extracted by identifying contiguous intervals with $I > 0.01$\,A.
Voltage and current are downsampled to \SI{1}{\second} resolution via linear interpolation.
The reference SOC is obtained by Coulomb counting:
$\mathrm{SOC}(t) = 1 - \int_0^t I(\tau)\,\mathrm{d}\tau / Q_{\mathrm{meas}}$,
where $Q_{\mathrm{meas}}$ is the measured discharge capacity of the cycle.

\paragraph{Data split protocol}
Batch-1 Cycle-1 data is used for parameter fitting; Batch-1 Cycles 2--$n$ are used for same-batch validation.
Batch-2 data is reserved entirely for cross-batch testing.
No Batch-2 data is used for fitting or hyperparameter selection.

\subsection{Baseline models}\label{sec:baselines}
We compare the proposed Shepherd-based framework against four voltage models spanning the interpretability--complexity spectrum:
\begin{enumerate}[leftmargin=*, label=\arabic*.]
    \item \textbf{Nernst-Based Model (NBM):}
    A thermodynamically motivated OCV model augmented with an exponential correction:
    \begin{equation}\label{eq:nbm}
    V_{\mathrm{OC}} = E^0 + \frac{R_g T}{nF}\ln\!\frac{\mathrm{SOC}}{1-\mathrm{SOC}} + M\exp(-\alpha\cdot\mathrm{SOC}),
    \end{equation}
    with $V_{\mathrm{term}} = V_{\mathrm{OC}} - IR$.  Parameters: $E^0 = \SI{4.5}{\volt}$, $M = -1.1718$, $\alpha = 0.8022$, fitted by nonlinear least squares.
    \item \textbf{Rint (Internal Resistance):}
    $V = \mathrm{OCV}(\mathrm{SOC}) - R_{\mathrm{int}} I$, where OCV is a 7th-order polynomial fit and $R_{\mathrm{int}}$ is identified from the steady-state voltage--current characteristic.  Parameters: 9 (polynomial coefficients + $R_{\mathrm{int}}$).
    \item \textbf{Th\'{e}venin-1RC:}
    First-order ECM: $V = \mathrm{OCV} - R_0 I - V_{RC}$, $\dot{V}_{RC} = I/C_1 - V_{RC}/(R_1 C_1)$.  Parameters: 10 (polynomial OCV + $R_0, R_1, C_1$).
    \item \textbf{Th\'{e}venin-2RC:}
    Second-order ECM adding a diffusion RC branch: $V = \mathrm{OCV} - R_0 I - V_{RC,1} - V_{RC,2}$.  Parameters: 12.
\end{enumerate}
All ECM parameters are identified using differential-evolution global optimisation with identical settings (population 50, 1000 generations) and cost-function stabilisation (RC voltage clamping $|V_{RC}| \le \SI{2}{\volt}$ and finite-value guards) to ensure a fair comparison.

\paragraph{Note on data-driven baselines}
We do not include purely data-driven models (e.g., LSTM, XGBoost) in the quantitative voltage-tracking comparison because (i)~such models require representative training data from each target operating condition, whereas the proposed framework generalises from a single constant-current discharge curve; (ii)~data-driven models do not provide the component-level power attribution or interpretable temperature--aging coupling that underpins the energy management strategy (\secref{sec:energy_mgmt}); and (iii)~the XJTU dataset's constant-current protocol does not supply the diverse load profiles needed to train a competitive data-driven predictor.
An LSTM inference time is not included in the main comparison, as differences in hardware and implementation preclude fair benchmarking.

\begin{figure*}[!tbp]
\centering
\includegraphics[width=\textwidth]{fig_voltage_trajectory.pdf}
\caption{Voltage prediction comparison on XJTU Batch-1 Battery-1 (Cycle 0, 2\,C constant-current discharge). (a)~Terminal voltage trajectories for all five models and the experimental reference. (b)~Point-wise prediction error. The Shepherd model (red solid line) tracks the experimental curve with substantially lower error than the ECM and Nernst baselines throughout the full discharge.}\label{fig:voltage_trajectory}
\end{figure*}

\subsection{Evaluation metrics}\label{sec:metrics}
Three categories of metrics are employed.
First, \emph{voltage accuracy} is assessed via root-mean-square error (RMSE), mean absolute error (MAE), maximum absolute error (MaxErr), and mean relative error (MRE), computed point-wise over full discharge trajectories.
Second, \emph{TTE accuracy} is measured by the absolute and percentage deviation of ablated-model TTE from the full-model prediction.
Third, \emph{statistical robustness} is quantified through 95\% bootstrap confidence intervals obtained by Monte Carlo resampling: 500 parametric perturbation samples with $\pm5\%$ uniform noise on each parameter.

\subsection{Ablation protocol}\label{sec:ablation_protocol}
To quantify each submodel's contribution, five configurations are evaluated:
(i)~\emph{Full Model} with all submodels active;
(ii)~\emph{w/o Temperature}, fixing $T = \SI{25}{\degreeCelsius}$ ($R^{\mathrm{th}} = R_0$, $Q_0^{\mathrm{th}} = Q_0^{\mathrm{nom}}$);
(iii)~\emph{w/o Aging}, setting $N = 0$, $t_s = 0$;
(iv)~\emph{w/o Component Decomposition}, replacing \eqnref{eq:ptotal} with flat average power;
and (v)~\emph{w/o Polarisation}, setting $K = 0$ in \eqnref{eq:shepherd}.
Each configuration is evaluated across a $3 \times 3$ grid: temperatures $T \in \{0, 25, 40\}$\,\si{\degreeCelsius} and aging levels $N \in \{0, 300, 800\}$ equivalent full cycles, under the Reading scenario.

\subsection{Usage scenarios}\label{sec:scenarios}
Five representative smartphone usage patterns are defined in \tabref{tab:scenarios}, spanning the power spectrum from idle standby to high-intensity online gaming.
A \SI{5000}{\milli\ampere\hour} cell with $T = \SI{25}{\degreeCelsius}$ and $N = 0$ serves as the reference configuration.
\begin{table*}[!tbp]
\caption{Component utilisation profiles ($\xi_i$) for five usage scenarios.}\label{tab:scenarios}
\centering
\begin{tabular}{@{}lccccc@{}}
\toprule
Component & Standby & Navigation & Gaming & Video Streaming & Reading \\
\midrule
Screen ($B$, $\eta$) & 0\%, 0 & 70\%, 1.0 & 90\%, 1.3 & 80\%, 1.0 & 70\%, 1.0 \\
SoC ($\xi_1$) & 10\% & 30\% & 100\% & 30\% & 30\% \\
Radio ($\xi_2$) & 10\% & 30\% & 95\% & 90\% & 20\% \\
GPS\textsuperscript{a} & On & On & On & On & On \\
Baseline ($\xi_3$) & 100\% & 100\% & 100\% & 100\% & 100\% \\
\midrule
Total Power (W) & 0.741 & 1.845 & 4.303 & 2.598 & 1.727 \\
\bottomrule
\multicolumn{6}{@{}p{0.95\textwidth}@{}}{\footnotesize \textsuperscript{a}GPS is modelled as binary on/off; $P_{\mathrm{GPS}} = \SI{300}{\milli\watt}$ when active. All scenarios assume background location services are active.}
\end{tabular}
\end{table*}

% ============================================================
\section{Results and Discussion}\label{sec:results}
\subsection{Model validation}\label{sec:validation}
\subsubsection{Overall voltage prediction accuracy}\label{sec:voltage}
\tabref{tab:baseline} compares point-wise voltage-tracking performance across all five models on XJTU Batch-1 (mean $\pm$ std over 8 cells).
\begin{table}[!tb]
\caption{Voltage prediction accuracy on XJTU Batch-1 (mean $\pm$ std, 8 cells).}\label{tab:baseline}
\centering
\resizebox{\columnwidth}{!}{\small
\begin{tabular}{@{}lcccc@{}}
\toprule
Model & RMSE (mV) & MAE (mV) & MaxErr (mV) & MRE (\%) \\
\midrule
NBM & $72.70 \pm 1.23$ & $49.99 \pm 0.58$ & 719.5 & 1.410 \\
Rint & $39.79 \pm 0.66$ & $16.88 \pm 0.28$ & 555.7 & 0.500 \\
Th\'{e}venin-1RC & $39.75 \pm 0.51$ & $17.02 \pm 0.24$ & 553.7 & 0.510 \\
Th\'{e}venin-2RC & $39.81 \pm 0.44$ & $17.33 \pm 0.50$ & 551.7 & 0.510 \\
\textbf{Shepherd (Proposed)} & $\mathbf{17.85 \pm 0.44}$ & $\mathbf{9.75 \pm 0.72}$ & \textbf{294.1} & \textbf{0.280} \\
\bottomrule
\end{tabular}
}
\end{table}
The proposed Shepherd model achieves an RMSE of \SI{17.85}{\milli\volt}, representing a \SI{75.5}{\percent} improvement over NBM and a \SI{55.1}{\percent} improvement over the best ECM baseline (Th\'{e}venin-1RC).
The improvement is consistent across all eight cells (coefficient of variation $< 3\%$), confirming that the accuracy gains are trajectory-wide rather than localised to specific SOC regions.
\figref{fig:voltage_trajectory} visualises the voltage predictions and point-wise errors for a representative cell.


\paragraph{Why the three ECMs perform identically}
The three ECM baselines yield nearly identical RMSE values (\SIrange{39.75}{39.81}{\milli\volt}).
This can be understood from the nature of the XJTU discharge data: constant-current (CC) protocols produce a monotonically declining voltage trajectory with no transient load steps.
Under CC conditions, the RC branch voltages in Th\'{e}venin models equilibrate within the first few seconds, after which the RC dynamics contribute a near-constant offset indistinguishable from adjusting $R_0$.
Consequently, the additional parameters in 1RC and 2RC models are structurally redundant for CC data and cannot improve upon the Rint baseline.

\paragraph{Source of the Shepherd advantage}
The Shepherd model's $K/\mathrm{SOC}$ polarisation term provides a \emph{nonlinear} voltage correction that increases as the cell depletes.
This captures the concentration-gradient-driven overpotential that intensifies at low SOC---a physical effect that constant-valued ECM resistances cannot represent.
The exponential transient zone ($A$, $B$) further accounts for the initial voltage sag upon load application, which the Nernst model's logarithmic singularity erroneously amplifies at SOC extremes (explaining the NBM's high MaxErr of \SI{719.5}{\milli\volt}).

\paragraph{Empirical justification of the polarisation exponent}
A natural extension of \eqnref{eq:shepherd} is to generalise the polarisation term from $K/\mathrm{SOC}$ to $K/\mathrm{SOC}^{\,p}$, treating the exponent $p$ as a sixth free parameter. Refitting the generalised six-parameter model on Batch-1 by differential evolution with $p \in [0.3, 2.5]$ yields $p^{*} = 1.026$, with a pooled voltage RMSE of \SI{17.65}{\milli\volt}---an improvement of less than \SI{0.02}{\milli\volt} over the classical $p = 1$ form (\SI{17.66}{\milli\volt}). This empirical near-equality of $p^{*}$ to unity, observed independently across all eight cells, confirms that the Shepherd polarisation kernel is already at the data-supported optimum and that introducing the additional parameter is statistically unjustified for NMC/graphite cells under constant-current loading. The result also explains why the classical formulation, despite its 1965 vintage, has remained competitive with more elaborate concentration-overpotential expressions for engineering use.

\paragraph{Why does $p^{*}$ land so close to one?}
The number $p^{*}=1.026$ is small enough to look like a fitting artefact, but it is not. There is a short electrochemical argument behind it. In the Shepherd ansatz the term $K/\mathrm{SOC}$ stands in for the concentration overpotential that builds up at the particle--electrolyte interface during discharge. If we linearise Butler--Volmer about open circuit and use Fick's first law across the diffusion layer, the leading-order overpotential scales as $\eta_\mathrm{conc} \propto i / c_\mathrm{s}$, with $c_\mathrm{s}$ the surface lithium concentration. Now stoichiometry: over the band $\mathrm{SOC}\in[0.1, 0.9]$, where the cells spend more than 95\% of their discharge time, $c_\mathrm{s}$ is, to a good approximation, linear in SOC. So the inverse-linear ($p=1$) form is what the physics predicts, and a clear departure $|p-1|\gtrsim 0.1$ would be diagnostic of something we are missing---either solid-state diffusion taking over (which would push $p$ towards $0.5$, the Warburg fingerprint), or a two-phase plateau being crossed (which would inflate $p$ above $1.5$). Neither happens here. The fitted $p^{*}=1.026 \pm 0.04$ across all eight Batch-1 cells therefore tells us something we did not assume going in: the NMC/graphite chemistry, under the C-rates studied, stays in the charge-transfer-limited regime where the original 1965 Shepherd ansatz is physically appropriate. The parsimony of the proposed model is then less a modelling concession than a statement about the cell.

\paragraph{On the choice of a classical formulation}
The Shepherd equation dates to 1965, yet its selection here is deliberate rather than retrograde.
More recent approaches---fractional-order ECMs, physics-informed neural networks, electrochemical impedance spectroscopy (EIS)-parameterised models---offer richer dynamics but at the cost of $>$20 parameters, specialised measurement equipment (EIS analysers), or GPU-accelerated training data that conflict with the portable-deployment objective of this work.
The present contribution is not a new voltage model \emph{per se}, but a \emph{system-level integration} that couples a right-sized voltage core with temperature, aging, and component-level load submodels to solve the discharge-side energy management problem.
We argue that model complexity should be \emph{commensurate with the application}: for the 1-second TTE update loop envisioned on an ARM Cortex-M microcontroller, the Shepherd formulation occupies the Pareto-optimal point of accuracy vs.\ computational cost, achieving lower RMSE than all ECM baselines with the fewest parameters (5 vs.\ 9--12).

\paragraph{A practical guideline for picking the right model}
Reading \tabref{tab:baseline} and \tabref{tab:nasa_val} side by side, a fairly clear decision rule falls out. If the load is quasi-steady or piecewise-constant---portable electronics, low-frequency grid storage, slow EV cruising---the proposed Shepherd model is the right pick: it sits on the Pareto front of accuracy versus parameter count, and nothing in our results suggests it would be worth the extra complexity to go elsewhere. The picture changes when the dominant frequency content of the load climbs above $\sim\!1/\tau_{RC} \approx \SI{0.1}{\hertz}$, which is the regime of urban EV drive cycles, frequency-regulation duty, and pulsed DC fast-charging. There the diffusion-limited transient becomes visible and a single Th\'{e}venin RC branch grafted onto the Shepherd voltage core recovers the missing dynamic fidelity while still keeping the five Shepherd parameters interpretable; we mention this Shepherd--RC hybrid as the natural next step, not as a finished result. Finally, for off-line lifetime simulation on a workstation, design-of-experiment studies, or anything else where neither parameter parsimony nor real-time execution is a hard constraint, the single-particle and pseudo-2D electrochemical models~\cite{Doyle1993,Mama2025} remain the reference and should be used. The present paper targets the first regime, and \secref{sec:dynamic_val} delineates its boundary by reporting honestly how the model degrades when pushed into the second.

\subsubsection{Cross-batch generalisation}\label{sec:crossbatch}
To test transferability, we apply the parameters fitted on Batch-1 (2\,C, 8 cells) directly to Batch-2 (3\,C, 15 cells) without any re-fitting.
The same-batch RMSE is $21.58 \pm 1.78$\,mV (MAE $14.71 \pm 1.98$\,mV, MaxErr 280.1\,mV), while the cross-batch result is $20.36 \pm 3.68$\,mV (MAE $12.96 \pm 4.71$\,mV, MaxErr 284.0\,mV).
The cross-batch RMSE is marginally \emph{lower} than the same-batch value, with the bootstrap-estimated gap of $\Delta = -1.22$\,mV (95\% CI: $[-3.88, +0.07]$\,mV) spanning zero.
This confirms that the model generalises across C-rates without accuracy degradation.
\figref{fig:cross_batch} shows the per-cell RMSE distributions for both conditions.
\begin{figure}[!htb]
\centering
\includegraphics[width=\columnwidth]{fig_cross_batch.pdf}
\caption{Per-cell RMSE distributions for same-batch (Batch-1, 2\,C) and cross-batch (Batch-2, 3\,C, no re-fitting) evaluation. Individual data points overlaid on violin plots. The distributions overlap substantially, confirming transferability.}\label{fig:cross_batch}
\end{figure}
\emph{Discussion.}
The slightly lower cross-batch RMSE is not paradoxical: Batch-2 cells (3\,C) exhibit a more compressive voltage profile that happens to be closer to the Shepherd model's implicit polarisation scaling at higher currents.
The key observation is that the confidence interval for the gap includes zero, indicating no statistically significant accuracy difference.

\begin{figure*}[!tbp]
\centering
\includegraphics[width=\textwidth]{fig_nasa_dynamic_validation.pdf}
\caption{Dynamic validation on NASA random-walk data (RW5, representative cycle). (a)~Measured and predicted terminal voltage under the random-walk current profile. (b)~Applied current sequence ($0.5$--$4$\,A). (c)~Point-wise voltage prediction error.}\label{fig:nasa_val}
\end{figure*}

\subsubsection{Validation under dynamic load profiles}\label{sec:dynamic_val}
The preceding subsections validate the Shepherd model under constant-current protocols.
To assess whether CC-identified parameters generalise to time-varying currents---a prerequisite for the dynamic TTE predictions in \secref{sec:tte_analysis}---we evaluate the model on the NASA Randomized Battery Usage Dataset~\cite{Bole2014}.
For each cell (RW3--RW6), the Shepherd parameters ($E_0$, $K$, $A$, $B$, $R_0$) are fitted on up to five reference CC discharges using L-BFGS-B minimisation of voltage RMSE, achieving CC fit RMSE $\le \SI{26}{\milli\volt}$.
The fitted model is then applied \emph{without re-identification} to 50 random-walk (RW) discharge cycles per cell (200 cycles total), where the current profile consists of consecutive \SI{5}{\minute} segments drawn uniformly from $[0.5, 4]$\,A.
Voltage is simulated via Euler integration of the SOC dynamics with the measured current profile; RMSE is computed over active discharge segments only (rest intervals excluded).

To benchmark the Shepherd transferability against a representative dynamic-capable baseline, we replicate the same CC-fit$\to$RW-test protocol with a Th\'{e}venin-1RC equivalent-circuit model. For each NASA cell, the OCV(SOC) polynomial and $(R_0, R_1, C_1)$ are jointly identified by differential evolution on the reference CC discharge (with bounds $R_0 \in [0.005, 0.3]$\,$\Omega$, $R_1 \in [0.001, 0.2]$\,$\Omega$, $C_1 \in [50, 5\!\times\!10^4]$\,F) and then propagated through each random-walk cycle by Euler-integrating the RC and SOC states under the measured $I(t)$. \tabref{tab:nasa_val} summarises the per-cell and aggregate results for both models.
\begin{table}[!tb]
\caption{Dynamic validation on NASA random-walk data: the proposed Shepherd model versus a Th\'{e}venin-1RC ECM, both calibrated only on the reference CC discharges and applied to the random-walk cycles \emph{without re-fitting}. Values are mean $\pm$ std across cycles; aggregate is over all 200 cycles.}\label{tab:nasa_val}
\centering
\resizebox{\columnwidth}{!}{\small
\begin{tabular}{@{}lccccc@{}}
\toprule
Cell & Model & CC fit RMSE (mV) & RW RMSE (mV) & RW MAE (mV) & Cycles \\
\midrule
\multirow{2}{*}{RW3} & Shepherd            & 22.3 & $86.9 \pm 15.9$  & $69.7 \pm 15.3$ & 50 \\
                     & Th\'{e}venin-1RC    & 19.3 & $183.9 \pm 34.1$ & $150.4 \pm 35.3$& 50 \\
\multirow{2}{*}{RW4} & Shepherd            & 25.6 & $42.1 \pm 12.1$  & $32.6 \pm 11.0$ & 50 \\
                     & Th\'{e}venin-1RC    & 18.4 & $202.4 \pm 34.0$ & $170.6 \pm 38.5$& 50 \\
\multirow{2}{*}{RW5} & Shepherd            & 24.3 & $37.6 \pm 11.5$  & $29.3 \pm 10.3$ & 50 \\
                     & Th\'{e}venin-1RC    & 18.2 & $199.6 \pm 37.2$ & $166.9 \pm 38.4$& 50 \\
\multirow{2}{*}{RW6} & Shepherd            & 19.6 & $87.3 \pm 16.6$  & $69.6 \pm 16.5$ & 50 \\
                     & Th\'{e}venin-1RC    & 14.8 & $214.1 \pm 36.3$ & $174.6 \pm 40.8$& 50 \\
\midrule
\multirow{2}{*}{All} & \textbf{Shepherd}            & ---  & $\mathbf{63.5 \pm 27.6}$  & $\mathbf{50.3 \pm 23.7}$ & 200 \\
                     & Th\'{e}venin-1RC             & ---  & $200.0 \pm 37.0$ & $165.6 \pm 38.5$ & 200 \\
\bottomrule
\end{tabular}
}
\end{table}
\figref{fig:nasa_val} shows a representative RW cycle from battery RW5, illustrating the model's ability to track the staircase-like voltage response to randomly switching currents.


\emph{Discussion.}
The aggregate RW RMSE of \SI{63.5}{\milli\volt} represents a $\sim 3.5\times$ increase over the CC fit RMSE ($\sim$\SI{23}{\milli\volt}), reflecting the additional challenges posed by dynamic currents: instantaneous ohmic drops at current transitions ($I \cdot \Delta R_0$), diffusion-limited transients not captured by the quasi-static Shepherd formulation, and cumulative SOC tracking errors under high C-rate segments.
The cell-to-cell variation (RW4/RW5 at $\sim$\SI{40}{\milli\volt} vs.\ RW3/RW6 at $\sim$\SI{87}{\milli\volt}) likely reflects differences in cell health and impedance characteristics.
Nonetheless, a sub-\SI{100}{\milli\volt} RMSE across 200 dynamic cycles---achieved with parameters identified solely from CC data on a \emph{different cell chemistry and format} than the XJTU cells---demonstrates that the Shepherd formulation provides a reasonable first-order approximation even under strongly time-varying loads.
\emph{Comparison with a dynamic-capable baseline.} Strikingly, the head-to-head benchmark in \tabref{tab:nasa_val} shows that the Th\'{e}venin-1RC ECM, despite carrying an extra dynamic state $V_{RC}$, achieves an aggregate RW RMSE of $200.0 \pm 37.0$\,mV---roughly $3.2\times$ worse than the Shepherd model under the identical CC-fit$\to$RW-test protocol. The reason mirrors the structural-identifiability discussion of \secref{sec:paramfit}: the Th\'{e}venin RC time constants $(R_1, C_1)$ are essentially unidentifiable from CC data because no current step exists to excite the exponential RC mode, so the optimiser converges to a parameter set whose RC dynamics are effectively dormant during fitting but are then exposed---catastrophically---by the random-walk current transitions. The Shepherd model, by contrast, has no dynamic state to mis-identify: its polarisation and exponential terms are functions of $(I, \mathrm{SOC})$ alone and therefore carry over cleanly from CC to RW. This observation reinforces the model-selection guideline of \secref{sec:voltage}: ECMs should be parameterised on data containing genuine current transitions (pulse tests, RW segments, or hybrid pulse--CC protocols), not on CC discharges alone.
For applications requiring higher fidelity than the present Shepherd-only formulation under dynamic conditions, a Th\'{e}venin RC branch fitted on \emph{dynamic} reference data, or a neural residual corrector, could be incorporated as a modular extension (\secref{sec:conclusion}).

\paragraph{How does the model fare against generic data-driven surrogates?}
It is fair to ask how a five-parameter physics-grounded model holds up against the obvious modern alternative: just train a black-box voltage map. We did this on the same NASA cells with two surrogates of increasing capacity---a degree-3 polynomial ridge regression (10 parameters) and a tiny multilayer perceptron with one hidden layer of 5 tanh units, trained by L-BFGS (21 parameters). To make the comparison as unfavourable to our model as possible, we let both surrogates see, on top of the reference CC discharges, the first ten random-walk cycles of each cell as a calibration set. Evaluation was then performed on the held-out RW cycles 11--50, which neither model had seen. \tabref{tab:datadriven} collects the test-set RMSE.
\begin{table}[!tb]
\centering
\caption{Data-driven surrogates compared with the proposed Shepherd model on held-out NASA random-walk cycles 11--50. The data-driven baselines were additionally calibrated on RW cycles 1--10 of each cell; the Shepherd model saw only the reference CC discharges. Mean $\pm$ std across $4 \times 40 = 160$ cycles.}\label{tab:datadriven}
{\footnotesize
\begin{tabular}{@{}lrrr@{}}
\toprule
Model & \#~params & RMSE (mV) & RW data needed \\
\midrule
Shepherd (proposed)        &  5 & $62.2 \pm 28.2$ & none \\
Polynomial-3 ridge surrogate & 10 & $28.0 \pm \phantom{0}6.1$ & 10 calib.\ cycles per cell \\
MLP (2--5--1) surrogate    & 21 & $26.5 \pm \phantom{0}6.2$ & 10 calib.\ cycles per cell \\
\bottomrule
\end{tabular}
}
\end{table}
Once they are allowed to see in-domain RW data, both surrogates do beat Shepherd on raw RMSE, by a factor of about $2.4\times$. This is what universal approximators are supposed to do, and we are not surprised. The cost of that gain, though, comes in three forms. First, parameter count: $2$--$4\times$ more than Shepherd. Second, the calibration requirement---ten labelled RW cycles per cell, which a deployed BMS often simply does not have. Third, and to us the most consequential, the loss of any handle on temperature, aging, or component decomposition of the kind that \secref{sec:model} couples through physical parameters. A surrogate would need to relearn all of that, separately, for every new $(T, N, \text{usage})$ combination it meets in the field. Read in this light, the $62$\,mV is not a defeat. It is a different operating point: only the Shepherd model generalises from CC data alone---five numbers per cell, obtainable from a single end-of-line characterisation discharge---to arbitrary random-walk loads with no retraining. And in the embedded scenario we are after (\secref{sec:comp_cost}), neither polynomial expansion nor neural inference fits the Cortex-M execution budget for sub-second updates over the lifetime of the cell. We pay $\sim\!35$\,mV of additional voltage RMSE for parsimony, zero-calibration transfer, and physical interpretability. We think the trade is worth making.

\begin{figure*}[!tbp]
\centering
\includegraphics[width=\textwidth]{fig_residual_analysis.pdf}
\caption{Residual diagnostics for the Shepherd model (Batch-1, Battery-1). (a)~Voltage error vs.\ time, coloured by SOC (\%). (b)~Error histogram with fitted normal overlay and Shapiro--Wilk test. (c)~SOC-segmented RMSE and MAE with mean bias (red diamonds, right axis).}\label{fig:residual_analysis}
\end{figure*}

\subsubsection{Bootstrap confidence intervals}\label{sec:bootstrap}
Bootstrap 95\% CIs (500 parametric perturbation samples, $\pm5\%$ uniform noise on each parameter) confirm statistical reliability: Shepherd RMSE $= 21.58$\,mV [19.99, 22.38]; TTE ranges from $3.80$\,h [3.67, 3.94] (Gaming) to $22.65$\,h [21.87, 23.51] (Standby); cross-batch gap $= -1.22$\,mV [$-3.88$, $+0.07$].
All CIs are narrow (typically $< 8\%$ of the point estimate).

\subsubsection{Residual analysis}\label{sec:residual}
To evaluate the quality of the Shepherd model's voltage predictions beyond aggregate metrics, we perform a detailed residual analysis on Batch-1, Battery-1.
\figref{fig:residual_analysis} presents three diagnostic views of the prediction residuals.
The SOC-coloured error scatter (panel~a) reveals that the largest deviations occur at SOC extremes---slight overestimation near SOC~$\approx 1$ (early discharge transient) and underestimation near SOC~$\approx 0.1$ (deep-discharge regime)---while mid-SOC errors remain tightly bounded within $\pm 10$\,mV.
The error histogram (panel~b) closely follows a fitted normal distribution ($\mu \approx 0$, indicating negligible systematic bias), with a Shapiro--Wilk test confirming approximate normality and suggesting that the residual errors are predominantly stochastic.
The SOC-segmented analysis (panel~c) decomposes RMSE and MAE across four SOC ranges, quantifying the expected increase in prediction difficulty at SOC extremes and providing actionable guidance for targeted model refinement.

\subsection{TTE prediction and scenario analysis}\label{sec:tte_analysis}
\subsubsection{Steady-state usage scenarios}\label{sec:tte}
\tabref{tab:tte} reports predicted TTE values at $T = \SI{25}{\degreeCelsius}$, $N = 0$ cycles, and $Q_0 = \SI{5.0}{\ampere\hour}$.
\begin{table}[!tb]
\caption{Predicted TTE and total power for each scenario (\SI{5000}{\milli\ampere\hour} cell, \SI{25}{\degreeCelsius}, new battery).}\label{tab:tte}
\centering
\begin{tabular}{@{}lcc@{}}
\toprule
Scenario & Total Power (W) & TTE (h) \\
\midrule
Standby & 0.741 & 22.64 \\
Navigation & 1.845 & 9.02 \\
Gaming (Online) & 4.303 & 3.79 \\
Video Streaming & 2.598 & 6.37 \\
Reading & 1.727 & 9.64 \\
\bottomrule
\end{tabular}
\end{table}
\figref{fig:scenario_comparison} visualises TTE and component-level power for all five scenarios.

Online gaming depletes the battery $5.9\times$ faster than standby, driven by simultaneous high utilisation of SoC (100\%), screen at 120\,Hz (90\% brightness), and cellular radio (95\%).
Navigation's TTE (\SI{9.02}{\hour}) is close to Reading (\SI{9.64}{\hour}) despite notably different load compositions: Navigation draws more radio power ($\xi_2 = 30\%$ vs.\ 20\%) but a similar SoC and screen profile, yielding only a \SI{6.8}{\percent} difference in total power (\SI{1.845}{\watt} vs.\ \SI{1.727}{\watt}) and hence comparable runtimes.

\begin{figure*}[!tbp]
\centering
\includegraphics[width=\textwidth]{fig_scenario_comparison.pdf}
\caption{(a)~Time-to-empty for five usage scenarios (\SI{5000}{\milli\ampere\hour} cell, \SI{25}{\degreeCelsius}, new battery). Bar labels show TTE (top) and total power (centre). (b)~Stacked component power decomposition revealing which subsystem dominates each scenario.}\label{fig:scenario_comparison}
\end{figure*}

\paragraph{Component power breakdown}
For Video Streaming, the radio subsystem dominates at 1055.7\,mW (40.6\%), followed by SoC at 612.0\,mW (23.6\%), screen at 510.6\,mW (19.7\%), GPS at 300.0\,mW (11.5\%), and the baseline at 120.0\,mW (4.6\%), consistent with sustained data-streaming demands.

\subsubsection{Dynamic mixed-usage scenario}\label{sec:dynamic}
To demonstrate the framework's applicability under realistic, time-varying workloads, we simulate a sequential multi-activity session:
Phase~1: 30\,min online gaming (\SI{4.303}{\watt}),
Phase~2: 60\,min video streaming (\SI{2.598}{\watt}),
Phase~3: 90\,min reading (\SI{1.727}{\watt}),
Phase~4: standby (\SI{0.741}{\watt}) until depletion.
\figref{fig:dynamic_scenario} shows the resulting voltage and SOC trajectories.
The total predicted TTE is \SI{15.7}{\hour}---significantly longer than constant gaming (\SI{3.79}{\hour}) but shorter than constant standby (\SI{22.6}{\hour})---illustrating how the framework can compose heterogeneous usage phases via piecewise-constant power scheduling.
The sharp SOC slope during the gaming phase (Phase~1) and gradual flattening during standby (Phase~4) are clearly visible, demonstrating the model's ability to capture the distinct energy drain rates of different workloads within a single simulation.


\subsection{Sensitivity, ablation, and multi-condition analysis}\label{sec:multi_analysis}
\subsubsection{Ablation study}\label{sec:ablation}
\figref{fig:ablation_heatmap} displays the resulting TTE error heatmaps, and \tabref{tab:ablation} reports the corresponding numerical values across a $3 \times 3$ temperature--aging grid under the Reading scenario.


\begin{table}[!tb]
\caption{Ablation study: TTE error (\%) introduced by removing each submodel (Reading scenario).}\label{tab:ablation}
\centering
{\footnotesize
\begin{tabular}{@{}llrrrr@{}}
\toprule
Temp.\ & Aging & w/o Temp.\ & w/o Aging & w/o Comp.\ & w/o Polar.\ \\
       &       & (\%)       & (\%)      & (\%)       & (\%)       \\
\midrule
\multirow{3}{*}{\SI{0}{\degreeCelsius}}
& New (0 cyc.) & $+70.6$ & $0.0$ & $+21.8$ & $+0.9$ \\
& 300 cyc. & $+70.2$ & $+13.6$ & $+21.8$ & $+0.8$ \\
& 800 cyc. & $+69.6$ & $+53.4$ & $+21.9$ & $+0.8$ \\
\midrule
\multirow{3}{*}{\SI{25}{\degreeCelsius}}
& New (0 cyc.) & $0.0$ & $0.0$ & $+21.8$ & $+0.9$ \\
& 300 cyc. & $0.0$ & $+13.9$ & $+21.9$ & $+0.9$ \\
& 800 cyc. & $0.0$ & $+54.3$ & $+22.0$ & $+0.9$ \\
\midrule
\multirow{3}{*}{\SI{40}{\degreeCelsius}}
& New (0 cyc.) & $-18.5$ & $0.0$ & $+21.8$ & $+0.9$ \\
& 300 cyc. & $-18.4$ & $+14.0$ & $+21.9$ & $+0.9$ \\
& 800 cyc. & $-18.2$ & $+54.8$ & $+22.0$ & $+0.9$ \\
\bottomrule
\end{tabular}
}
\end{table}

\begin{figure*}[!tbp]
\centering
\includegraphics[width=\textwidth]{fig_dynamic_scenario.pdf}
\caption{Dynamic mixed-usage scenario: 30\,min Gaming $\to$ 60\,min Video $\to$ 90\,min Reading $\to$ Standby until depletion. Top panel: terminal voltage trajectory with phase-coloured background shading. Bottom panel: SOC evolution. The different slopes reflect the distinct power demands of each usage phase.}\label{fig:dynamic_scenario}
\end{figure*}

\begin{figure*}[!tbp]
\centering
\includegraphics[width=\textwidth]{fig_ablation_heatmap.pdf}
\caption{Ablation heatmaps showing TTE error (\%) when each submodel is removed, across a $3\times3$ temperature--aging grid (Reading scenario). (a)~Without temperature coupling: up to $+70\%$ error at \SI{0}{\degreeCelsius}. (b)~Without aging: up to $+55\%$ at 800 cycles. (c)~Without component decomposition: $\approx +22\%$ (up to $+43\%$ for Gaming). (d)~Without polarisation ($K=0$): $<1\%$.}\label{fig:ablation_heatmap}
\end{figure*}

\paragraph{Temperature coupling}
This is the dominant submodel at off-nominal temperatures: ablation at \SI{0}{\degreeCelsius} causes a $+70\%$ TTE overestimate (the model falsely assumes room-temperature capacity), while ablation at \SI{40}{\degreeCelsius} causes a $-18\%$ underestimate (missing the capacity gain from increased lattice occupancy, partially offset by leakage).
The asymmetry ($|+70\%| \gg |-18\%|$) reflects the nonlinear logistic shape of $Q_0(T)$, which is steeper below $T_m$ than above.

\paragraph{Aging dynamics}
Aging importance grows monotonically with cycle count---from exactly 0\% at $N=0$ (by construction) to $\sim$55\% at 800 cycles---and is nearly independent of temperature ($\pm 1\%$ variation across the three $T$ levels).
This near-separability validates the multiplicative coupling assumed in \eqnref{eq:Q0coupled}.

\paragraph{Component decomposition}
Replacing the five-component model with flat average power induces a consistent $\approx +22\%$ TTE error for the Reading scenario across all temperature and aging conditions.
The maximum error reaches $+43\%$ under the Gaming scenario, where the divergence between the detailed five-component and the lumped power model is largest due to the high dynamic range of concurrent SoC, screen, and radio loads.

\paragraph{Polarisation}
The $K$ term contributes only $0.8$--$0.9\%$ to TTE error, consistent with its primary role in mid-trajectory \emph{voltage shape} rather than integrated \emph{energy capacity}.
Nevertheless, removing $K$ visibly worsens voltage RMSE and would degrade SOC-estimation accuracy in real-time BMS applications.

\paragraph{Non-additive coupling justifies the unified framework}
Summing the individual ablation errors for, e.g., \SI{0}{\degreeCelsius} and 800 cycles ($70.6 + 53.4 + 21.9 + 0.8 = 146.7\%$) far exceeds the actual combined error observed in any single ablation, confirming that the submodel contributions are not independent.
The multiplicative coupling in \eqnsref{eq:Q0coupled}{eq:Rcoupled} produces interaction terms---e.g., temperature-reduced capacity amplifies aging-induced resistance---that a simple additive combination of standalone submodels would miss.
This non-additivity is precisely why a \emph{unified} framework is needed rather than a pipeline of independent corrections.

\subsubsection{Parameter sensitivity analysis}\label{sec:sensitivity}
A local one-at-a-time sensitivity analysis is conducted by perturbing each parameter by $\pm 5\%$ around the Video Streaming baseline (TTE $= \SI{6.37}{\hour}$), following the rapid parameter estimation and sensitivity methodology advocated by V\'{a}zquez-Arenas et al.~\cite{VazquezArenas2014ECM}.
\begin{table}[!tb]
\caption{Parameter sensitivity ($\pm 5\%$ perturbation, Video Streaming).}\label{tab:sensitivity}
\centering
\resizebox{\columnwidth}{!}{\small
\begin{tabular}{@{}lcccc@{}}
\toprule
Parameter & $t_{-5\%}$ (h) & $t_{+5\%}$ (h) & $\Delta_{-5\%}$ (\%) & $\Delta_{+5\%}$ (\%) \\
\midrule
$P_{\mathrm{total}}$ & 6.71 & 6.06 & $+5.37$ & $-4.86$ \\
$Q_0$ & 6.06 & 6.67 & $-4.77$ & $+4.76$ \\
$E_0$ & 6.05 & 6.68 & $-4.96$ & $+4.96$ \\
$K$ & 6.37 & 6.36 & $+0.07$ & $-0.07$ \\
$R_0$ & 6.37 & 6.37 & $+0.04$ & $-0.04$ \\
$A$ & 6.35 & 6.38 & $-0.24$ & $+0.24$ \\
$B$ & 6.38 & 6.35 & $+0.24$ & $-0.22$ \\
\bottomrule
\end{tabular}
}
\end{table}
TTE is primarily governed by three ``energy-budget'' parameters---$P_{\mathrm{total}}$, $Q_0$, and $E_0$ (sensitivity $\approx \pm 5\%$)---whose product $Q_0 E_0 / P_{\mathrm{total}}$ approximates TTE to first order.
The internal Shepherd parameters ($K$, $R_0$, $A$, $B$) contribute less than \SI{0.3}{\percent} each, confirming that they primarily shape the \emph{voltage trajectory} rather than the \emph{integrated runtime}.
This hierarchy has practical implications: for TTE estimation, accurate calibration of capacity and load power is far more critical than precise identification of internal model parameters.
\figref{fig:sensitivity_tornado} presents the corresponding tornado chart.


\subsubsection{Temperature impact on TTE}\label{sec:temp_tte}
Under the Video Streaming scenario, TTE varies from 3.44\,h at $-10\,^\circ$C to a peak of $\sim$6.37\,h near $25$--$35\,^\circ$C, before declining to 5.82\,h at $50\,^\circ$C.
The unimodal pattern arises from the competition between two temperature-dependent effects.
Below \SI{35}{\degreeCelsius}, reduced ionic conductivity and contracted lattice occupancy decrease effective capacity and increase resistance, compressing TTE; the effect is severe, with TTE at $-10\,^\circ$C reaching only 53\% of the peak value.
Above \SI{35}{\degreeCelsius}, although capacity continues to increase modestly, rapidly rising SoC leakage power (\eqnrefp{eq:leak}; incremental $+\SI{0.35}{\watt}$ at \SI{35}{\degreeCelsius} vs.\ $+\SI{1.22}{\watt}$ at \SI{50}{\degreeCelsius} relative to the \SI{25}{\degreeCelsius} reference) more than offsets the capacity gain, causing TTE to decline.
This non-monotonic behaviour is consistent with field observations that smartphones shut down prematurely in both extreme cold and extreme heat.
\figref{fig:temperature_tte} illustrates the complete temperature--TTE response.
\begin{figure}[!htb]
\centering
\includegraphics[width=\columnwidth]{fig_temperature_tte.pdf}
\caption{TTE vs.\ ambient temperature for the Video Streaming scenario (\SI{5000}{\milli\ampere\hour}, new battery). The unimodal curve peaks near \SI{35}{\degreeCelsius}; performance degrades sharply in cold environments (capacity loss) and moderately in hot environments (leakage rise).}\label{fig:temperature_tte}
\end{figure}


\subsubsection{Aging impact and degradation trajectories}\label{sec:aging_results}
\figref{fig:aging_curves} visualises the aging model's capacity fade and resistance growth trajectories alongside the resulting TTE impact.
Under the Video Streaming scenario, TTE degrades from \SI{6.4}{\hour} (new battery) to approximately \SI{3.3}{\hour} after 1000 equivalent full cycles, a reduction of $\sim$48\%.
The super-linear growth exponent ($\gamma = 1.085$) indicates accelerating degradation, consistent with progressive SEI thickening and pore clogging mechanisms reported in the aging literature~\cite{Vetter2005}.





\subsection{Engineering implications and energy management}\label{sec:eng_mgmt}
\subsubsection{Engineering implications}\label{sec:engineering}
The results yield several actionable insights for power management:
\begin{enumerate}[leftmargin=*, label=(\roman*)]
    \item \textbf{Load-aware scheduling.}
    Since gaming depletes the battery $5.9\times$ faster than standby, OS-level awareness of load composition enables proactive low-battery warnings and adaptive quality-of-service adjustments.
    The component decomposition framework provides the necessary attribution for identifying which subsystem to throttle.
    \item \textbf{Temperature-aware operation.}
    The $\sim 70\%$ TTE overestimation at \SI{0}{\degreeCelsius} when temperature is ignored motivates incorporating ambient temperature into the OS battery estimator---a feature absent from most current mobile BMS implementations.
    \item \textbf{Health-aware budgeting.}
    After 800 cycles, TTE drops by $\sim 55\%$.
    Proactive notification and automatic power-profile switching for aged batteries can prevent unexpected shutdowns.
    \item \textbf{SOC management.}
    The calendar aging model (\eqnrefp{eq:calaging}) quantitatively shows that maintaining SOC within the 20\%--80\% range extends battery lifespan by reducing parasitic SEI growth at high electrode potentials.
\end{enumerate}



\subsubsection{Adaptive energy management optimisation}\label{sec:energy_mgmt}
To demonstrate the framework's utility as a \emph{model-predictive control} (MPC) prediction engine for discharge-side power scheduling, we formulate the runtime-aware throttling problem as a receding-horizon optimisation and compare it against a rule-based threshold baseline.
This subsection illustrates one downstream application of the validated Shepherd model; the modelling contributions of \secref{sec:model}--\secref{sec:results} stand independently of any specific control strategy.

\paragraph{MPC formulation}
Let $T_{\mathrm{des}}$ denote the user's desired runtime (e.g., $T_{\mathrm{des}} = \SI{6}{\hour}$).
At each control instant $t_k$ (period $\Delta t_c = \SI{120}{\second}$), the MPC solves:

\begin{equation}\label{eq:mpc}
\begin{aligned}
\min_{\mathbf{u}_k} \; J_k = \sum_{j=0}^{N_p - 1} \Bigl[
& \underbrace{\alpha \Bigl(\frac{P_{\max} - u_{k+j}}{P_{\max}}\Bigr)^{\!2}}_{\text{QoS loss}}
+ \underbrace{\beta \bigl(\hat{s}_{k+j} - s_{k+j}^{\mathrm{ref}}\bigr)^{2}}_{\text{SOC tracking}} \\
& + \underbrace{\gamma \Bigl(\frac{u_{k+j} - u_{k+j-1}}{P_{\max}}\Bigr)^{\!2}}_{\text{smoothness}}
\Bigr]
\end{aligned}
\end{equation}
subject to:
\begin{align}
\hat{s}_{k+j+1} &= \hat{s}_{k+j} - \frac{u_{k+j}\,\Delta t_c}{3600\,V_{\mathrm{Shep}}(\hat{s}_{k+j}, T, N) \cdot Q_0(T,N)}, \label{eq:mpc_dyn} \\
P_{\min} &\le u_{k+j} \le P_{\max}, \quad j = 0,\ldots,N_p{-}1, \label{eq:mpc_bounds}
\end{align}

\begin{figure}[!htb]
\centering
\includegraphics[width=\columnwidth]{fig_sensitivity_tornado.pdf}
\caption{Tornado chart of parameter sensitivity ($\pm 5\%$ perturbation, Video Streaming scenario). TTE is dominated by three energy-budget parameters ($P_{\mathrm{total}}$, $Q_0$, $E_0$); internal Shepherd parameters contribute $<0.3\%$ each.}\label{fig:sensitivity_tornado}
\end{figure}

where $\mathbf{u}_k = [u_k,\ldots,u_{k+N_p-1}]^\top$ is the power decision vector over a prediction horizon of $N_p = 10$ steps ($\SI{1200}{\second}$), $\hat{s}$ denotes the predicted SOC, $V_{\mathrm{Shep}}$ is the Shepherd voltage model (\eqnrefp{eq:shepherd}), and $s^{\mathrm{ref}}(t) = s(t_k) - [s(t_k) - s_{\mathrm{cut}}] \cdot (t - t_k)/(T_{\mathrm{des}} - t_k)$ is a linear reference trajectory from the current SOC to the cutoff threshold $s_{\mathrm{cut}} = 0.05$ at $t = T_{\mathrm{des}}$.
Only the first action $u_k^*$ is applied; the horizon then shifts forward and the optimisation is re-solved---the standard receding-horizon principle.
The weights $(\alpha, \beta, \gamma) = (0.3, 500, 0.2)$ are selected to prioritise SOC tracking (ensuring the battery reaches $T_{\mathrm{des}}$) while maximising quality of service.
\eqnref{eq:mpc} is solved via L-BFGS-B with box constraints (\eqnrefp{eq:mpc_bounds}), requiring $< 60$ iterations per step.

\paragraph{Threshold baseline}
For comparison, a two-stage rule-based strategy applies fixed power reductions at predetermined SOC levels: Stage~1 (SOC~$\le 20\%$) reduces power from \SI{4.303}{\watt} (gaming) to \SI{1.73}{\watt} (reading-level); Stage~2 (SOC~$\le 10\%$) further reduces to \SI{0.74}{\watt} (standby).

\paragraph{Scenario}
A user begins online gaming ($P_{\max} = \SI{4.303}{\watt}$) at full charge with $T_{\mathrm{des}} = \SI{6}{\hour}$.
The baseline (no management) maintains constant gaming power until depletion at \SI{3.8}{\hour}.

\paragraph{Results}
\tabref{tab:mgmt_results} and \figref{fig:energy_mgmt} compare the three strategies.
\begin{table}[!tb]
\caption{Energy management strategy comparison (Gaming, \SI{5000}{\milli\ampere\hour}, \SI{25}{\degreeCelsius}).}\label{tab:mgmt_results}
\centering
\resizebox{\columnwidth}{!}{\small
\begin{tabular}{@{}lcccc@{}}
\toprule
Strategy & TTE (h) & Extension (\%) & Energy (Wh) & Efficiency gain (\%) \\
\midrule
No management & 3.79 & --- & 16.3 & --- \\
Threshold (rule-based) & 5.32 & $+40$ & 16.4 & $+39$ \\
\textbf{MPC ($T_{\mathrm{des}} = \SI{6}{\hour}$)} & \textbf{5.99} & $\mathbf{+58}$ & \textbf{16.5} & $\mathbf{+56}$ \\
\bottomrule
\end{tabular}
}
\end{table}



The MPC strategy achieves a TTE of \SI{5.99}{\hour}---within \SI{0.2}{\percent} of the target---extending runtime by \SI{58}{\percent} relative to the baseline.
This exceeds the threshold strategy (\SI{5.32}{\hour}, $+40\%$) by an additional 18 percentage points because the MPC \emph{anticipates} the energy budget:
it begins gradually reducing power well before the SOC thresholds, producing a smooth, monotonically decreasing power profile (\hyperref[fig:energy_mgmt]{\textup{Figure~\ref*{fig:energy_mgmt}b}}) rather than the abrupt step changes of the rule-based approach.
The SOC trajectory (\hyperref[fig:energy_mgmt]{\textup{Figure~\ref*{fig:energy_mgmt}a}}) closely tracks the linear reference budget, confirming that the receding-horizon controller successfully regulates energy consumption to meet the target runtime.

\begin{figure*}[!tbp]
\centering
\includegraphics[width=\textwidth]{fig_aging_curves.pdf}
\caption{(a)~Capacity fade $Q_0(N)$ and internal resistance growth $R_0(N)$ as functions of equivalent full cycles. (b)~Corresponding TTE degradation under the Video Streaming scenario, with percentage loss annotations.}\label{fig:aging_curves}
\end{figure*}

\paragraph{Discussion}
The MPC formulation offers three structural advantages over threshold strategies.
First, the power trajectory is \emph{continuous} rather than piecewise-constant, enabling smoother quality-of-service degradation that is less perceptible to the user.
Second, the \emph{target runtime is an explicit input} ($T_{\mathrm{des}}$), allowing the OS to expose a ``last until'' setting that adapts the throttling profile automatically---unlike fixed-threshold approaches that are insensitive to the user's remaining schedule.
Third, because the prediction model (\eqnrefp{eq:mpc_dyn}) inherits the temperature and aging coupling of the Shepherd framework, the MPC adapts its power trajectory to the battery's current health: an aged cell ($N = 800$) or a cold environment ($T = \SI{0}{\degreeCelsius}$) triggers earlier and deeper throttling \emph{without any parameter re-tuning}, whereas rule-based thresholds remain static.

\paragraph{Connection to prior battery-management work}
The MPC strategy represents a \emph{discharge-side} counterpart to the economy-conscious charging optimisation studied by Liu et al.~\cite{Liu2019ECM,Liu2018ECM} in the EV domain.
Their coupled electrothermal-aging models jointly optimise charging time, temperature rise, and degradation cost; the present work extends this model-predictive philosophy to the discharge side, using the Shepherd model as the internal prediction engine for load reduction.
The component-level power decomposition (\eqnrefp{eq:ptotal}) identifies \emph{which} subsystem to throttle---a capability that mirrors the component-level thermal attribution in battery thermal management systems~\cite{Mama2025}.
At the aggregate scale, the framework's ability to predict TTE under arbitrary $(T, N)$ conditions could support coordinated charge scheduling across large device fleets, contributing to peak-shaving and demand-response objectives in smart-grid scenarios.

\paragraph{Energy efficiency quantification}
Over the baseline depletion window (\SI{3.8}{\hour}), the MPC consumes approximately \SI{14.2}{\watt\hour} vs.\ the baseline's \SI{16.3}{\watt\hour}, yielding a per-cycle energy reduction of $\sim$13\% at matched service duration.
If the full MPC session (\SI{5.99}{\hour}) is considered, the energy efficiency---defined as useful runtime per unit energy consumed---improves by approximately \SI{56}{\percent}:
\begin{equation}\label{eq:mpc_eff}
\eta_{\mathrm{MPC}} = \frac{t_{\mathrm{MPC}} / E_{\mathrm{MPC}}}{t_{\mathrm{base}} / E_{\mathrm{base}}} = \frac{5.99 / 16.5}{3.79 / 16.3} \approx 1.56.
\end{equation}
Note that these figures are simulated values under idealised load assumptions; real-world gains will depend on load-profile accuracy and user acceptance of quality-of-service degradation.

\paragraph{Comparison with commercial strategies}
Existing commercial power-saving modes---such as Android Battery Saver, iOS Low Power Mode, and Samsung Adaptive Power Saving---typically apply fixed or coarsely staged power reductions that are neither target-runtime-aware nor responsive to battery health state; their internal implementations are proprietary and not publicly benchmarked under standardised conditions.
By contrast, the proposed MPC accepts an explicit target runtime as input, incorporates the battery's current temperature and aging state into the optimisation via the Shepherd prediction model, and produces a continuous power trajectory rather than abrupt step changes.
These structural differences are complementary rather than competitive: the MPC logic could in principle be integrated into existing OS-level power managers to improve their adaptiveness.
The computational overhead is negligible: each MPC solve requires $< 60$ L-BFGS-B iterations over 10 decision variables, completing in $< \SI{1}{\milli\second}$ on a single ARM core, well within the \SI{120}{\second} control interval.

\begin{figure*}[!tbp]
\centering
\includegraphics[width=\textwidth]{fig_energy_management.pdf}
\caption{Energy management comparison: no management (dashed grey), threshold (dash-dot teal), and MPC with $T_{\mathrm{des}} = \SI{6}{\hour}$ (solid red). (a)~SOC trajectories; the dotted grey line shows the linear SOC budget. The MPC trajectory tracks the budget closely, reaching TTE~$= \SI{6.0}{\hour}$ ($+58\%$).  (b)~Power profiles: the MPC produces a smooth, gradually decreasing power schedule vs.\ the threshold's abrupt step changes.  (c)~Cumulative energy delivered over time.}\label{fig:energy_mgmt}
\end{figure*} 


% ============================================================
\section{Conclusions and Future Work}\label{sec:conclusion}
This paper presented a unified continuous-time simulation framework for portable-device battery TTE prediction that integrates a modified Shepherd voltage model, five-component power decomposition, logistic capacity--temperature and Arrhenius resistance--temperature coupling, and power-law cycle/calendar aging into a single differential-algebraic system.
The Shepherd-based formulation achieves a voltage RMSE of \SI{17.85 \pm 0.44}{\milli\volt}, reducing error by \SI{75.5}{\percent} relative to the Nernst baseline and \SI{55}{\percent} relative to Th\'{e}venin-1RC on XJTU Batch-1.
Cross-batch validation (Batch-1$\to$Batch-2, no re-fitting) confirms generalisability across C-rates, with a statistically non-significant RMSE gap ($-1.22$\,mV, 95\% CI spanning zero).
Dynamic-load validation on the NASA random-walk dataset~\cite{Bole2014} yields a voltage RMSE of \SI{63.5 \pm 27.6}{\milli\volt} across 200 cycles with currents spanning 0.5--\SI{4}{\ampere}, using parameters identified solely from constant-current data, demonstrating that the Shepherd formulation transfers to time-varying loads without re-identification.

Multi-condition ablation across a $3 \times 3$ temperature--aging grid quantifies each module's contribution: temperature coupling up to \SI{70}{\percent}, aging up to \SI{55}{\percent}, component decomposition up to \SI{43}{\percent}, and polarisation $<\SI{1}{\percent}$ of TTE error.
A receding-horizon MPC energy management strategy extends simulated runtime by up to \SI{58}{\percent} (TTE $= \SI{6.0}{\hour}$ vs.\ $\SI{3.8}{\hour}$ baseline) and improves energy efficiency by \SI{56}{\percent}, outperforming the threshold-based strategy ($+40\%$) by an additional 18 percentage points while producing a smooth, continuous power profile.
The framework's computational cost of \SI{0.009}{\milli\second} per simulated second confirms suitability for real-time deployment on embedded hardware.

\paragraph{Scope and limitations}
This work is positioned as a \emph{simulation framework} validated on cell-level electrochemical data.
The voltage model and its parameters are fitted and validated on real cell data---both under constant-current conditions (XJTU dataset) and under dynamic random-walk loads (NASA dataset~\cite{Bole2014})---while the component-level load profiles are synthetic, constructed from published hardware power measurements~\cite{Carroll2010,Chen2014,Lu2008}.
Although the TTE predictions under these synthetic loads cannot be directly verified against real-device measurements, the dual-dataset voltage validation (CC and dynamic) ensures a sound electrochemical foundation.
Importantly, the synthetic load profiles are not arbitrary---they are constructed by combining individually measured component powers from peer-reviewed hardware studies, and they are applied as inputs to an independently validated electrochemical model.
This decoupled validation approach (electrochemistry on real cell data, loads from literature) is common in the battery modelling literature~\cite{Fuller2014,Mama2025} and provides a rigorous basis for scenario exploration.
The framework's modular load interface (\ssref{sec:power}) is designed as a plug-and-play architecture: replacing the synthetic profiles with real telemetry data (e.g., from Android's BatteryStats API or iOS's Energy Diagnostics) requires only substituting the $\xi_i(t)$ time series, with no changes to the electrochemical or thermal submodels.
Additional limitations are noted below.
The model assumes a constant ambient temperature per simulation; rapid thermal transients from SoC self-heating are not captured.
As emphasised by Mama et al.~\cite{Mama2025}, electrochemical-thermal coupled (ECTC) models are essential for accurately predicting temperature-dependent degradation under sustained high-power loads.
Liu et al.~\cite{Liu2019ECM} demonstrated that even a lumped two-state thermal model significantly improves fidelity; Jeon and Baek~\cite{Jeon2011ECM} and Liu et al.~\cite{Liu2017ECM} further established thermal modelling methodologies for lithium-ion cells during discharge. Incorporating a similar thermal ODE is a natural next step.
Moreover, the algebraic Shepherd formulation does not capture RC-type transient relaxation dynamics; incorporating a Th\'{e}venin RC branch or re-identifying parameters on standardised dynamic profiles (e.g., DST, UDDS) could further reduce the dynamic RMSE.
The present model treats nominal capacity $Q_0$ as a function of temperature and aging only (\eqnrefp{eq:Q0T}) without a Peukert-style correction for reduced effective capacity at high discharge rates~\cite{Fuller2014}; for portable-device cells operating predominantly between \SIrange{0.3}{1.5}{C}, this omission has limited practical impact.
Finally, cell-to-cell variability and electrode-level spatial heterogeneity are not modelled, and results are on NMC/graphite cells---other chemistries (LFP, NCA, solid-state) require re-parameterisation.

\paragraph{Future work}
Planned extensions include:
(i)~end-to-end validation using real device power traces with concurrent battery telemetry, building on the dynamic-load validation presented in \secref{sec:dynamic_val};
(ii)~coupled electro-thermal modelling via a lumped thermal ODE for self-heating;
(iii)~stochastic load modelling using Markov-chain usage profiles;
(iv)~integration into Android's PowerManagerService or Linux power governance for on-device field testing;
(v)~hybrid physics-data correction using a lightweight neural residual network to compensate for unmodelled transient dynamics (e.g., the RC relaxation effects observed in the NASA dynamic validation);
and (vi)~exploration of aggregated distributed portable storage for vehicle-to-grid (V2G) style demand response, where billions of portable devices collectively contribute to grid frequency regulation and peak shaving through coordinated charge scheduling.

% ============================================================
\section*{CRediT authorship contribution statement}
\textbf{Yinuo Wang}: Methodology, Software, Validation, Formal analysis, Data curation, Visualisation, Writing -- original draft.
\textbf{Yike Yang}: Conceptualisation, Methodology, Writing -- review \& editing.
\textbf{Lizhi Wei}: Conceptualisation, Methodology, Writing -- review \& editing.
\textbf{Zhiqiang Li}: Supervision, Project administration, Writing -- review \& editing.

\section*{Declaration of competing interest}
The authors declare that they have no known competing financial interests or personal relationships that could have appeared to influence the work reported in this paper.

\section*{Data availability}
The XJTU battery dataset is publicly available~\cite{XJTU2020}. The NASA Randomized Battery Usage Dataset is available from the NASA Prognostics Data Repository~\cite{Bole2014}. Code and processed data supporting this study are available at \url{https://github.com/wangyinuo001/BatteryPaper}.

\section*{Acknowledgements}
The authors gratefully acknowledge the XJTU research group for making their battery dataset publicly available, and the NASA Prognostics Center of Excellence for providing the Randomized Battery Usage Dataset.

% ============================================================
\appendix
\section{Proof of \propref{prop:mono}}\label{app:proof}
\begin{proof}
For fixed SOC, define $F = E_0 + A\exp\!\bigl(-BQ_0(1-\mathrm{SOC})\bigr) > 0$ and $\alpha = K/\mathrm{SOC} + R > 0$.
The Shepherd voltage (\eqnrefp{eq:shepherd}) is affine in~$I$: $V_{\mathrm{term}}(I) = F - \alpha I$.
Hence $g(I) = I\,V_{\mathrm{term}}(I) = FI - \alpha I^2$ is a concave quadratic with $g(0) = 0$ and a unique maximum $g^* = F^2/(4\alpha)$ at $I^* = F/(2\alpha)$.
For any $P_{\mathrm{total}} \in (0, g^*]$, the equation $g(I) = P_{\mathrm{total}}$ has discriminant $\Delta = F^2 - 4\alpha P_{\mathrm{total}} \ge 0$, yielding two roots $I_{\pm} = (F \pm \sqrt{\Delta})/(2\alpha)$.
The smaller root $I_- \le I^*$ gives $V_{\mathrm{term}}(I_-) = F - \alpha I_- \ge F/2 > 0$, a physically valid operating point.
Since $g'(I) = F - 2\alpha I > 0$ for $I < I^*$, $g$ is strictly increasing on $[0, I^*]$ and $I_-$ is the unique root in this interval.
\end{proof}

%% Bibliography
\bibliographystyle{model1-num-names}
\bibliography{references}


\end{document}
